\documentclass[11pt]{article}
\usepackage[a4paper,margin=26mm]{geometry}
\usepackage[T1]{fontenc}
\usepackage{lmodern,amsmath,amssymb,amsthm,mathtools,booktabs,microtype}
\usepackage[colorlinks=true,linkcolor=blue,citecolor=blue,urlcolor=blue]{hyperref}
\newtheorem{theorem}{Theorem}[section]
\newtheorem{proposition}[theorem]{Proposition}
\newtheorem{lemma}[theorem]{Lemma}
\newtheorem{corollary}[theorem]{Corollary}
\theoremstyle{definition}\newtheorem{definition}[theorem]{Definition}
\theoremstyle{remark}\newtheorem{remark}[theorem]{Remark}
\newcommand{\E}{\mathbb E}\newcommand{\Pp}{\mathbb P}
\newcommand{\T}{\mathbb T}\newcommand{\R}{\mathbb R}\newcommand{\C}{\mathbb C}
\newcommand{\N}{\mathbb N}\newcommand{\Z}{\mathbb Z}
\newcommand{\dd}{\,\mathrm d}\newcommand{\ind}{\mathbf1}
\newcommand{\Var}{\operatorname{Var}}\newcommand{\Poi}{\operatorname{Poi}}
\newcommand{\Leb}{\operatorname{Leb}}\newcommand{\TV}{\operatorname{TV}}
\newcommand{\supp}{\operatorname{supp}}\newcommand{\Ree}{\operatorname{Re}}
\newcommand{\Op}{O_{\Pp}}\newcommand{\eps}{\varepsilon}
\numberwithin{equation}{section}\allowdisplaybreaks
\setlength{\emergencystretch}{2em}
\title{Spectral extremes under exact cycle conditioning}
\author{Zhipeng Lu\thanks{Shenzhen MSU--BIT University and Guangdong
Laboratory of Machine Perception and Intelligent Computing,
Shenzhen 518172, Guangdong, China.
Email: \href{mailto:zhipeng.lu@hotmail.com}{zhipeng.lu@hotmail.com}.}}
\date{September 10, 2026}
\hypersetup{pdftitle={Spectral extremes under exact cycle conditioning}}
\begin{document}\maketitle
\begin{abstract}
Let $P_n$ be the matrix of a random permutation of $n$ symbols and let
$M_n=\log\max_{|z|=1}|\det(I-zP_n)|$. Cook and Zeitouni proved that
$M_n/\log n$ converges in probability to a constant $x_0$ for a uniform
permutation. We show that the $\sqrt{\log n}$ fluctuations of $M_n$ are
carried entirely by the number of cycles $K_n$. Write
$\lambda(s)=\log\{\Gamma(1+s)/\Gamma(1+s/2)^2\}$, let $s_\kappa$ be the
minimizer of $(1+\kappa\lambda(s))/s$ on $(0,\infty)$, and put
$v(\kappa)=\kappa\lambda'(s_\kappa)$ and
$a_\theta=\lambda(s_\theta)/s_\theta$. Under the Ewens measure with any
fixed parameter $\theta>0$,
\[
 M_n=v(\theta)\log n+a_\theta\bigl(K_n-\theta\log n\bigr)
       +\Op(\log\log n),
\]
so that the standardized pair $(K_n,M_n)$ converges jointly to $(G,G)$
with $G$ standard normal: the maximum and the cycle count are
asymptotically perfectly aligned. This is deduced from a statement about
the exact conditional law, which does not depend on $\theta$: for every
compact $[\kappa_-,\kappa_+]\subset(0,\infty)$ there is a finite $C$ with
\[
 \sup_{\kappa_-\log n\le k\le\kappa_+\log n}
 \Pp\bigl(|M_n-v(k/\log n)\log n|>C\log\log n\ \big|\ K_n=k\bigr)
 \longrightarrow0,
\]
uniformly over exact and possibly atypical cycle counts. The proof keeps
the size and the cycle count simultaneously in a two-variable
coefficient extraction. Cycles longer than $n/(\log n)^4$ are reserved as
an analytic factor whose coefficients are flat under every size shift
produced by the shorter cycles; positivity then converts a scalar
coefficient asymptotic into a relative comparison of the entire
path-constrained measure, with an error that does not degrade with the
number of constraints or with the rarity of the event. The constrained
lower bound comes from pointwise saddle estimates for killed
convolutions along a dyadic chain of endpoint boxes.
\end{abstract}

\section{Introduction}\label{sec:intro}

\subsection{The maximum of a random permutation characteristic
polynomial}

Let $\sigma_n$ be a random permutation of $\{1,\dots,n\}$, let $P_n$ be
its permutation matrix, and let $C_j=C_j(\sigma_n)$ be its number of
$j$-cycles. Write $e(t)=\exp(2\pi it)$ and $\T=\R/\Z$ with normalized
Lebesgue measure. The characteristic polynomial factors along cycles,
\[
 F_n(z)=\det(I-zP_n)=\prod_{j\ge1}(1-z^j)^{C_j},
\]
so that on the unit circle
\begin{equation}\label{eq:field}
 \log|F_n(e(t))|=\sum_{j\ge1}C_j\log|1-e(jt)| .
\end{equation}
We study
\[
 M_n=\log\|F_n\|_\infty=\max_{t\in\T}\log|F_n(e(t))|,
\]
which is also the logarithm of the maximum modulus of $\det(zI-P_n)$.

The field \eqref{eq:field} is of log-correlated type. For a uniform
permutation the counts $C_j$ are asymptotically independent Poisson
variables of mean $1/j$, so \eqref{eq:field} is a sum over $\asymp\log n$
effective dyadic scales of frequencies, its variance is of order
$\log n$, and the covariance of two points decays like the logarithm of
their distance. Its increments are, however, far from Gaussian: the
summand $\log|1-e(jt)|$ is bounded above by $\log 2$ and has an
exponential left tail. The extremal behaviour of $M_n$ is therefore
governed by a genuine Legendre transform and not by a Gaussian rate. The
relevant cumulant generating function is that of $\log|1-e(U)|$ for a
uniform $U\in\T$,
\begin{equation}\label{eq:A}
 A(s)=\int_\T|1-e(t)|^s\dd t=\frac{\Gamma(1+s)}{\Gamma(1+s/2)^2},
 \qquad\lambda(s)=\log A(s),\qquad s>-1.
\end{equation}
For complex $z$ with $\Ree z>-1$ we write $A(z)$ for the same integral,
with $|1-e(t)|^z=\exp(z\log|1-e(t)|)$ off the roots and value zero at the
roots; it is analytic and equals the same gamma ratio.

Cook and Zeitouni \cite[Theorem~1.2]{CZ} proved that for a uniform
permutation $M_n/\log n$ converges in probability to a constant $x_0$,
characterized by $\lambda^*(x_0)=1$ for the Legendre transform
$\lambda^*$ of $\lambda$. Their proof combines arithmetic separation of
frequencies, a division into small and large arcs, two-point estimates
and a truncated second moment. Those mechanisms reappear below, at the
level of the constrained kernels; what is new here is the way the exact
combinatorial constraints are carried through them.

\subsection{Why the cycle count should be conditioned on}
\label{subsec:why}

The randomness in \eqref{eq:field} has two layers of very different
character. The arithmetic layer is the choice of the frequencies $j$
that occur; it produces the log-correlated structure and all the
difficulty in locating the maximum. The other layer is the total number
of cycles
\[
 K_n=\sum_{j\ge1}C_j,
\]
a single global statistic that merely counts how many summands
\eqref{eq:field} has. Under the Ewens measure with parameter $\theta>0$
one has $K_n=\theta\log n+\Op(\sqrt{\log n})$, and for the maximum this
fluctuation is not negligible. The extremal speed depends on the number
of summands, so a displacement of order $\sqrt{\log n}$ in $K_n$ moves
$M_n$ by the same order; and $\sqrt{\log n}$ dominates every correction
of size $\log\log n$. The leading fluctuation of $M_n$ is therefore
inherited from the count layer, not from the arithmetic one.

This suggests separating the two layers, and in the Ewens family the
separation is exact. The cycle counts satisfy
\[
 \Pp_{\theta,n}(C_j=c_j)=h_n(\theta)^{-1}
       \prod_{j\ge1}\frac{(\theta/j)^{c_j}}{c_j!},
 \qquad\sum_jjc_j=n,
\]
where
\begin{equation}\label{eq:basic-coeff}
 h_n(u)=[z^n](1-z)^{-u},\qquad a_{n,k}=[u^k]h_n(u),
 \qquad h_0(u)=1,
\end{equation}
and coefficients with an impossible index are zero. Conditioning on
$K_n=k$ cancels the factor $\theta^k$ and leaves
\begin{equation}\label{eq:conditional-law}
 \Pp_{n,k}(C_j=c_j)=a_{n,k}^{-1}
       \prod_j\frac{j^{-c_j}}{c_j!},
 \qquad\sum_jjc_j=n,\quad\sum_jc_j=k.
\end{equation}
This law does not depend on $\theta$. It is a microcanonical law, with
the size and the number of cycles both fixed exactly, and it is the
canonical object that remains once the count layer has been removed.

Working with \eqref{eq:conditional-law} is strictly more demanding than
working with the Ewens law itself. The two are related by the mixture
identity
\begin{equation}\label{eq:mixture}
 \Pp_{\theta,n}(A_n)=\sum_k\Pp_{\theta,n}(K_n=k)\,\Pp_{n,k}(A_n),
\end{equation}
and convergence of its left-hand side to zero gives no uniform bound on
the conditional probabilities on the right: at the mode
$\Pp_{\theta,n}(K_n=k)$ is only of order $(\log n)^{-1/2}$, and it is
exponentially smaller for counts away from $\theta\log n$. A statement
about the exact conditional law has to be proved directly, and every
estimate that carries the size constraint has to keep a \emph{relative}
error throughout. Producing such an estimate for a rare path event,
\emph{after} all of the constraints defining that event have been
imposed, is the technical heart of this paper.

Exact cycle conditioning also arises on its own. In \cite{Lu} the author
showed that for two uniformly random perfect matchings on $2n$ labels
the profile of alternating components is Ewens$(1/2)$; there the number
of components is the natural quantity to fix, and
Section~\ref{sec:consequences} records the resulting spectral statement.

\subsection{Results}\label{subsec:results}

For each $\kappa>0$ the function $s\mapsto(1+\kappa\lambda(s))/s$ has a
unique minimizer $s_\kappa$ on $(0,\infty)$, characterized by
\begin{equation}\label{eq:critical}
 \kappa\{s_\kappa\lambda'(s_\kappa)-\lambda(s_\kappa)\}=1.
\end{equation}
Indeed $\lambda''(s)>0$ is the variance of $\log|1-e(U)|$ under the
exponential tilt of order $s$, so the left-hand side vanishes at $s=0$,
has strictly positive derivative $\kappa s\lambda''(s)$ for $s>0$, and
tends to infinity by Stirling's formula. Put
$v(\kappa)=\kappa\lambda'(s_\kappa)$. Then, with $\kappa=k/\log n$,
\begin{equation}\label{eq:center}
 m_{n,k}:=\inf_{s>0}\frac{\log n+k\lambda(s)}{s}
        =\frac{\log n+k\lambda(s_\kappa)}{s_\kappa}
        =v(\kappa)\log n.
\end{equation}

\begin{theorem}[Exact-cycle localization]\label{thm:main}
Let $0<\kappa_-<\kappa_+<\infty$. There is a finite constant $C$,
depending only on this interval, such that
\begin{equation}\label{eq:main}
 \sup_{\substack{k\in\Z\\\kappa_-\log n\le k\le\kappa_+\log n}}
 \Pp_{n,k}\bigl(|M_n-m_{n,k}|>C\log\log n\bigr)\longrightarrow0.
\end{equation}
\end{theorem}

The constant is chosen before the error tolerance and serves every
admissible integer $k$ at once, including counts far from the typical
Ewens value. Both the upper and the lower half of \eqref{eq:main} are
proved under the exact constraint; neither is obtained by averaging over
$K_n$.

Theorem~\ref{thm:main} transfers back to the Ewens family. Fix
$\theta>0$ and put $a_\theta=\lambda(s_\theta)/s_\theta>0$. Since
$K_n/\log n\to\theta$ in probability, one may apply \eqref{eq:main} with
$\kappa=K_n/\log n$ and expand $v$ to first order at $\theta$; the
quadratic remainder is $\Op(1)$. With probability tending to one and
with a fixed constant,
\begin{equation}\label{eq:random-center}
 \bigl|M_n-v(\theta)\log n-a_\theta(K_n-\theta\log n)\bigr|
   \le C\log\log n,
\end{equation}
and consequently, by the classical central limit theorem for $K_n$,
\begin{equation}\label{eq:joint}
 \left(\frac{K_n-\theta\log n}{\sqrt{\theta\log n}},\
       \frac{M_n-v(\theta)\log n}{a_\theta\sqrt{\theta\log n}}\right)
 \Longrightarrow(G,G),\qquad G\sim\mathcal N(0,1).
\end{equation}
This is Corollary~\ref{cor:ewens}. The extremal speed is set by an
environment --- by how many cycles the permutation happens to have ---
and that environment is itself asymptotically Gaussian; the maximum
follows it. Statement \eqref{eq:joint} is joint weak convergence, and no
assertion about moments or about a limit for the Pearson correlation is
made.

\subsection{Discussion}\label{subsec:discussion}

\paragraph{The speed function.} The map $\kappa\mapsto v(\kappa)$ is
strictly increasing and strictly concave: implicit differentiation of
\eqref{eq:critical} gives \eqref{eq:v-derivatives}, namely
$v'(\kappa)=\lambda(s_\kappa)/s_\kappa>0$ and
$v''(\kappa)=-\{\kappa^3s_\kappa^3\lambda''(s_\kappa)\}^{-1}<0$. More
cycles raise the extremal speed, with diminishing returns. The first
derivative is exactly the coefficient $a_\theta$ in
\eqref{eq:random-center}, which is why the random centring is linear in
$K_n$ at the accuracy considered here. Since $\lambda'(s)<\log2$ for
every $s>0$ we have $v(\kappa)<\kappa\log2$, in agreement with the
deterministic bound $M_n\le K_n\log2$; the two are asymptotically equal
as $\kappa\downarrow0$, where a permutation with few cycles can make all
of its factors nearly extremal at once. At $\kappa=1$ equation
\eqref{eq:critical} is the Legendre-transform equation
$\lambda^*(v(1))=1$, so $v(1)=x_0$ is the Cook--Zeitouni constant and
Theorem~\ref{thm:main} contains their leading order.

\paragraph{The localization window.} The width $C\log\log n$ in
\eqref{eq:main} is the natural limit of the present method. For
branching random walks the second-order term of the maximum is
logarithmic in the number of generations, and the number of effective
generations here is of order $\log n$; a correction of size $\log\log n$
is therefore expected in \eqref{eq:main}, and Theorem~\ref{thm:main}
does not resolve it. We identify neither the coefficient of such a
correction, nor a limit law for the residual, nor an extremal point
process. On the other hand $\log\log n\ll\sqrt{\log n}$, so the window
is far below the scale of \eqref{eq:random-center} and \eqref{eq:joint};
those two statements are sharp at leading order.

\begin{remark}\label{rem:family}
Since \eqref{eq:conditional-law} does not involve $\theta$,
Theorem~\ref{thm:main} applies verbatim to every law on cycle profiles
of the form
$\Pp(c)\propto w\bigl(\sum_jc_j\bigr)\prod_jj^{-c_j}/c_j!$ with
$w\ge0$: conditioning any such law on its cycle count returns
\eqref{eq:conditional-law}. Whenever the induced count concentrates on a
compact set of values of $k/\log n$, it follows that
$|M_n-m_{n,K_n}|\le C\log\log n$ with probability tending to one for
that law. The Ewens family is the case $w(k)=\theta^k$.
\end{remark}

\subsection{Relation to previous work}\label{subsec:previous}

Beyond \cite{CZ}, several strands bear on the problem. Hughes,
Najnudel, Nikeghbali and Zeindler \cite{HNNZ} study moments,
multiplicative class functions and linear statistics of characteristic
polynomials under generalized Ewens measures, with generating functions
in cycle counts. Dang and Zeindler \cite{DZ} obtain joint central limit
theorems at finitely many fixed observation points under arithmetic
hypotheses on those points. Fran\c{c}ois \cite[Theorem~2.2]{Francois}
proves convergence of generalized-Ewens characteristic polynomials in
the topology of local uniform convergence on the open unit disk. Neither
a finite collection of observation points nor convergence in the
interior controls the maximum on the boundary, which is a supremum over
a family of constraints whose size grows with $n$. The central limit
theorem for $K_n$ used in \eqref{eq:joint} is classical; sharper uniform
asymptotics for the Ewens count distribution are in \cite{KMS}, and the
coefficient asymptotics behind \eqref{eq:basic-coeff} belong to the
classical theory of Stirling numbers of the first kind.

Random centring of an extremum by an environment has precedents. For
branching random walks in a time-inhomogeneous random environment,
Mallein and Mi\l o\'s \cite{MM} identify an environment-dependent centre
together with a logarithmic correction. The mechanism behind
\eqref{eq:random-center} is of that type, with the cycle count playing
the role of the environment, although the model here is a finite
permutation and the proof is self-contained: it uses elementary count
couplings, path estimates and second moments rather than an extremal
theorem for random environments.

Cook and Gu have announced related work on maxima of Poissonian
log-correlated fields \cite{CG}. Public abstracts from January 2025,
October 2025 and January 2026 describe refined maximal behaviour for
related random series or trigonometric polynomials, modelled on a
branching random walk in a random time-dependent environment, and the
research listing describes the work as in preparation. A complete
theorem statement was not available in the sources consulted on
September 10, 2026, so no comparison with a full theorem, and no claim
of priority over that work, is made here. We note only that the object
treated below is the exact finite permutation law under a precise
cycle-count constraint, for which \eqref{eq:mixture} shows that an
unconditional statement is not sufficient.

The coefficient viewpoint continues that of \cite{Lu}. The scalar
principles behind the reservoir estimate of Section~\ref{sec:reservoir}
originate in the hybrid method of Flajolet, Fusy, Gourdon, Panario and
Pouyanne \cite{FFGPP}; the estimate needed here has a cutoff that grows
with $n$ and a complex marker, and is proved directly. The contribution
of the coefficient step is the uniformity of its relative error after
all restrictions of a rare path event have been imposed on a positive
measure.

\subsection{Outline of the proof}\label{subsec:outline}

Throughout the proof $L=\log n$, $\ell=\log L$ and $b=\lfloor
n/L^4\rfloor$. The following is the skeleton; the constants are fixed in
the order set out in Section~\ref{sec:maximum}.

\paragraph{An analytic reservoir.} A fixed-observation hybrid estimate
cannot control the growing family of path restrictions that a maximum
requires. Instead we reserve every cycle longer than $b$. What remains
is a polynomial whose coefficients are positive measures recording the
heights $\log|1-e(jt)|$ at the observation points, and whose degree is
$O(Lb)=o(n)$ under the exact count constraint. The reserved factor
$R_b(z,u)=\exp\bigl(u\sum_{j>b}z^j/j\bigr)$ has one algebraic
singularity and an entire correction; a uniform Hankel estimate and a
saddle point in the marker $u$ show that its coefficients are flat under
every size shift the short cycles can produce, with relative error
$O(1/\ell)$. Because the comparison is between positive measures, it can
be applied \emph{after} all path projections and yields a relative
comparison of the whole restricted measure. Neither the number of
projections, nor the dimension of the space, nor the mass of the event
enters the error. This is Theorem~\ref{thm:positive-transfer}, and it is
what replaces a total-variation approximation whose absolute error would
exceed the probability of the event being measured.

\paragraph{A constrained kernel lower bound.} After the transfer, the
short cycles may be treated as independent harmonic samples in
logarithmic blocks. Grouping the blocks dyadically from both ends
produces $B\asymp\ell$ coarse groups of widths $\mathsf H_j$ that vary
by a bounded factor. Along a chain of endpoint boxes of width
$\sqrt{\mathsf H_j}$ we bound the killed convolution kernels from below
pointwise, using a bridge estimate for the tilted log-sine law
(Lemma~\ref{lem:bridge}). Integrating over the intermediate boxes
cancels the density factors $\mathsf H_j^{-1/2}$ against the box widths,
so that the surviving cost is $\sum_jE_j^2$ rather than
$B\max_jE_j^2$, where $E_j$ measures the deviation of the realized
counts in group $j$. This is what keeps the total cost polynomial,
$p_L\ge L^{-C_*}$ (Lemma~\ref{lem:boxes}), and it is what limits the
final window to $O(\ell)$.

\paragraph{Fourier comparison and a second moment.} On arithmetically
separated points the tilted block kernels are close, in total variation
$O(L^{-J})$ for any fixed $J$, to independent copies of the reference
convolution (Lemma~\ref{lem:smooth}). Since $J$ may be chosen after
$C_*$, that error is $o(p_L^2)$ and the comparison is valid
\emph{relative} to the rare event, not merely in absolute terms. A first
and second moment for the constrained functional then give a high point
in any prescribed set of angles of measure at least one half; the
matching upper bound follows from an integrated moment estimate whose
input is an entropy gap $\kappa b_s<1$ for the block envelope
(Lemma~\ref{lem:entropy}).

\paragraph{Restoring the remaining cycles.} The low-frequency cycles are
restored on a random but admissible set of angles, using only that the
lower bound of the previous step is uniform over such sets. The long
cycles are reinserted by a deterministic stability estimate for the
supremum norm under multiplication by $\prod(1-z^{j_v})$
(Lemma~\ref{lem:insertion}), whose loss is controlled by an exact
conditional factorial-moment identity for
$\bigl(\sum_{j\le b}jC_j\bigr)\bigl(\sum_{l>b}C_l/l\bigr)$
(Lemma~\ref{lem:crossmass}). Both operations move the centre by $O(\ell)$
only.

\subsection{Organization}\label{subsec:organization}

Section~\ref{sec:reservoir} constructs the reservoir, proves the
event-uniform coefficient comparison and identifies the law of the
marked block counts. Section~\ref{sec:kernel} proves the pointwise
killed-convolution estimates, the regularity of the count environment
and the box-chain lower bound. Section~\ref{sec:fourier} supplies the
uniform one- and two-point Fourier comparisons. Section~\ref{sec:maximum}
fixes the constants, runs the moment computation, restores the low and
long cycles and completes the proof of Theorem~\ref{thm:main}.
Section~\ref{sec:consequences} derives Corollary~\ref{cor:ewens}, records
the two-matching application and delimits what is and is not asserted.

\section{An analytic reservoir and a positive coefficient transfer}
\label{sec:reservoir}
Throughout the proof set
\begin{equation}\label{eq:parameters}
 L=\log n,\quad \ell=\log L,\quad b=\lfloor n/L^4\rfloor,
 \quad H_b=\sum_{j=1}^b\frac1j,\quad T=L-H_b=4\ell-\gamma+o(1).
\end{equation}
All assertions are for sufficiently large $n$. Define
\begin{equation}\label{eq:reservoir}
 R_b(z,u)=\exp\left(u\sum_{j>b}\frac{z^j}{j}\right)
 =(1-z)^{-u}\exp\{-uH_b(z)\},\qquad
 H_b(z)=\sum_{j=1}^b\frac{z^j}{j},
\end{equation}
and $r_{N,l}^{(b)}=[z^Nu^l]R_b(z,u)$.

\subsection{A uniform scalar coefficient estimate}
\begin{lemma}[Reservoir coefficient estimate]\label{lem:reservoir}
Fix a compact set $\mathcal U\subset\C$. Uniformly for $u\in\mathcal U$,
$b\ge2$, and $N\ge4b$,
\begin{align}\label{eq:reservoir-scalar}
 [z^N]R_b(z,u)
 &=e^{-uH_b}h_N(u)+E_{N,b}(u),\\[-2mm]
 |E_{N,b}(u)|
 &\le C\frac bN e^{-\Ree(u)H_b}N^{\Ree(u)-1}
       +Cb^C e^{-N/(2b)}.\nonumber
\end{align}
The estimate is absolute, including at zeros of $1/\Gamma(u)$.
\end{lemma}
\begin{proof}
Continue \eqref{eq:reservoir} to the disk of radius $1+1/b$, slit along
$[1,1+1/b]$, using the branch of $\log(1-z)$ analytic at zero. Put
$\rho=1+1/b$. On $|z|=\rho$ the quantity
$H_b(z)+\log(1-z)$ is bounded by an absolute constant on either side
of the slit. To verify this, for $|\arg z|\le2/b$ use
$|H_b(z)-H_b|\le Cb|z-1|\le C$ and
$\log|1-z|=-\log b+O(1)$. Otherwise the identity
\[
 H_b(z)+\log(1-z)=-\int_0^z\frac{w^b}{1-w}\dd w
\]
on the radial segment gives a bound $C/(b|\arg z|)\le C$; here
$|1-te^{i\varphi}|\ge c(|1-t|+|\varphi|)$ for
$0\le t\le\rho$, $|\varphi|\le\pi$, after harmless changes of constants.
Thus $R_b$ is bounded on the outer circle, uniformly on $\mathcal U$.
The comparison function $e^{-uH_b}(1-z)^{-u}$ has at worst polynomial
size in $b$ there. Their outer-contour contributions to the coefficient
are at most $Cb^C\rho^{-N}\le Cb^C e^{-N/(2b)}$.

Deform Cauchy's contour into the outer circle, both sides of the slit,
and an indentation $|z-1|=1/N$. For $|w|\le1/b$,
\[
 |H_b(1+w)-H_b|\le Cb|w|,\qquad
 |e^{-u(H_b(1+w)-H_b)}-1|\le Cb|w|.
\]
On the indentation, $|z|^{-N-1}\le C$ and its length is $O(N^{-1})$.
The contribution of the difference of the two functions is at most
$Ce^{-aH_b}N^{a-1}(b/N)$, where $a=\Ree u$. On either side of the slit
write $z=1+t$, $1/N\le t\le1/b$. The remaining contribution is bounded by
\[
 Ce^{-aH_b}b\int_{1/N}^{1/b}t^{1-a}e^{-Nt/2}\dd t
 \le Ce^{-aH_b}N^{a-1}\frac bN.
\]
The last integral is uniform for $a$ in a compact interval; its lower
limit is $1/N$, so it poses no integrability problem when $a\ge2$.
Branch moduli contribute bounded factors on $\mathcal U$. This proves
\eqref{eq:reservoir-scalar}, without moving the indentation to zero.

% BEGIN EXPANSION reservoir_contour
\paragraph{Uniformity of the contour bounds.}
Here and below a compact complex marker set is enclosed in a disk
$|u|\le M$, with $M$ fixed. The estimate on the outer circle can be
checked without an infinite series outside its disk of convergence.
For $z=\rho e^{i\varphi}$ with $|\varphi|>2/b$, use the radial integral
above and $|1-te^{i\varphi}|\ge c|\varphi|$. Its absolute value is at most
\[
 \frac{C}{|\varphi|}\int_0^\rho t^b\dd t
 =\frac{C\rho^{b+1}}{(b+1)|\varphi|}
 \le \frac{C}{b|\varphi|}.
\]
For $|\varphi|\le2/b$, one has $|z-1|\asymp b^{-1}$ and
$|z-1|\le C/b$. The elementary identity
$z^j-1=(z-1)\sum_{v=0}^{j-1}z^v$ gives
$|H_b(z)-H_b|\le Cb|z-1|$. Since
$H_b=\log b+O(1)$ and the branch argument of $1-z$ is bounded, this
also bounds $H_b(z)+\log(1-z)$ in the small-angle sector.
These arguments apply to the two boundary values at the slit.

For clarity, the remaining power of $N$ in the bank integral is obtained
by the substitution $v=Nt$:
\[
 b\int_{1/N}^{1/b}t^{1-a}e^{-Nt/2}\dd t
 =bN^{a-2}\int_1^{N/b}v^{1-a}e^{-v/2}\dd v
 \le C_M bN^{a-2},\qquad |a|\le M.
\]
The last constant is finite because the lower endpoint is one and the
upper tail is exponential, uniformly in $a$. On the indentation,
$|(1-z)^{-u}|\le C_M N^a$, the analytic correction is $O_M(b/N)$,
and the arc length is $O(1/N)$. All terms are therefore absolute
estimates of the same difference of analytic functions. In particular
none was obtained by division by $h_N(u)$ or by $1/\Gamma(u)$; their
possible zeros cause no exception.
% END EXPANSION reservoir_contour

\end{proof}

\begin{lemma}[Marker extraction]\label{lem:marker}
Uniformly for $k/\log n$ in a positive compact interval,
\begin{equation}\label{eq:akn}
 a_{n,k}=\frac{L^k}{n k!\,\Gamma(k/L)}\{1+O(L^{-1})\}.
\end{equation}
Let $b$ be as in \eqref{eq:parameters}, $0\le d\le C_0 Lb$,
$N=n-d$, and $T_N=\log N-H_b$. Uniformly when $l/T_N$ belongs to a
positive compact interval,
\begin{equation}\label{eq:reservoir-marker}
 r_{N,l}^{(b)}
 =\frac{T_N^l}{N l!\,\Gamma(l/T_N)}
   \left\{1+O\left(T_N^{-1}+\sqrt{T_N}\,b/N\right)\right\}.
\end{equation}
Consequently, if $l/T$ belongs to a fixed positive compact interval,
\begin{equation}\label{eq:flatness}
 \sup_{0\le d\le C_0Lb}
 \left|\frac{r_{n-d,l}^{(b)}}{r_{n,l}^{(b)}}-1\right|
 \le \eps_n,\qquad \eps_n=O(\ell^{-1}+L^{-3})=o(1).
\end{equation}
\end{lemma}
\begin{proof}
For bounded $u$, the gamma-ratio formula, interpreted analytically at
nonpositive integers, gives
\[
 h_N(u)=N^{u-1}\left\{\frac1{\Gamma(u)}+O(N^{-1})\right\},
\]
with an absolute uniform error of modulus $CN^{\Ree u-2}$.
For an entire function $g$, $x\to\infty$, and $j/x$ in a positive
compact interval, the elementary marker saddle is
\begin{equation}\label{eq:elementary-saddle}
 [u^j]e^{xu}g(u)=\frac{x^j}{j!}\{g(j/x)+O(x^{-1})\}.
\end{equation}
For completeness, write $g(u)=\sum g_m u^m$ and divide the left-hand
side by $x^j/j!$. It becomes
$\sum_{m\le j}g_m(j)_m/x^m$. On a disk of radius greater than twice the
largest $j/x$, Cauchy's coefficient bound and
$|1-(j)_m/j^m|\le m(m-1)/(2j)$ give the error $O(j^{-1})$.
The omitted tail is exponentially small. This proves
\eqref{eq:elementary-saddle}; it also holds for a function analytic on
such a fixed disk.

Apply it with $g=1/\Gamma$ and $x=L$. The coefficient of the gamma-ratio
error on $|u|=k/L$ is at most
$Cn^{-2}(k/L)^{-k}e^k$, which is smaller than the main term by
$O(\sqrt L/n)$. This proves \eqref{eq:akn}.
For the reservoir, Lemma~\ref{lem:reservoir} gives the main function
$N^{-1}e^{T_Nu}/\Gamma(u)$. On $|u|=l/T_N$, Cauchy's bound for its error,
compared with Stirling's estimate for $T_N^l/l!$, is
$O(\sqrt{T_N}b/N)$ relatively. The outer-contour error is negligible:
$N/b\asymp L^4$ whereas $\log N\asymp L$, so even a fixed polynomial
in $N,b$ times $e^{-N/(2b)}$ tends to zero faster than any negative
power of $L$. The main-function saddle gives the additional $O(T_N^{-1})$.

Finally $T_N-T=\log(N/n)=O(d/n)$, and the logarithmic derivative in $N$
of $T_N^l/(N\Gamma(l/T_N))$ is $O(N^{-1})$ uniformly in the indicated
range. Equivalently, take the ratio of the explicit main terms.
Together with \eqref{eq:reservoir-marker}, this proves \eqref{eq:flatness}.

% BEGIN EXPANSION marker_flatness
\paragraph{The two normalizations in the marker argument.}
To detail the power-series estimate, put $t=j/x$ and choose a fixed
radius $R_0>2\sup t$. If $|g_m|\le M_0R_0^{-m}$, then
\[
 \left|\sum_{m=0}^j g_m\frac{(j)_m}{x^m}-g(t)\right|
 \le \frac{M_0}{2j}\sum_{m=0}^j m(m-1)(t/R_0)^m
       +M_0\sum_{m>j}(t/R_0)^m
 =O(x^{-1}).
\]
The bound for $(j)_m$ follows from
$1-\prod_{v=0}^{m-1}(1-v/j)\le\sum_{v=0}^{m-1}v/j$.
This calculation is valid for complex coefficients as well. For
$g=1/\Gamma$ its value at the positive saddle is bounded away from
zero on the prescribed interval. Thus the absolute $O(x^{-1})$
error in the braces becomes a relative error there.

In the reservoir application, let $t=l/T_N$. The Cauchy error coming
from the first term of \eqref{eq:reservoir-scalar} is at most
$C(b/N^2)t^{-l}e^{tT_N}$. The positive main coefficient has size
$N^{-1}(T_N^l/l!)g(t)$. Their ratio is at most
\[
 C\frac bN\,\frac{e^l l!}{l^l}
 \le C\frac bN\sqrt l=O\!\left(\sqrt{T_N}\,b/N\right).
\]
For flatness, keep the integer $l$ and the cutoff $b$ fixed while the
real size variable $N$ moves from $n-d$ to $n$. If
$F(N)=T_N^l/(N\Gamma(l/T_N))$, and
$\psi=\Gamma'/\Gamma$ on the positive real axis, then
\[
 \frac{\dd}{\dd N}\log F(N)
 =\frac1N\left\{-1+\frac l{T_N}
            +\frac l{T_N^2}\psi(l/T_N)\right\}.
\]
The braces are uniformly bounded. Also $d/n=O(L^{-3})$ and
$T_N\asymp T\asymp\ell$. Integrating the derivative bounds the ratio
of main terms by $1+O(d/n)$; the two relative saddle errors add
$O(T^{-1}+\sqrt T\,b/n)$. This proves the stated $\eps_n$ uniformly
over every allowed size shift with a single constant.
% END EXPANSION marker_flatness

\end{proof}

\subsection{One transfer for the entire constrained kernel}
For finite positive measures, write $\mu\le\nu$ if the inequality holds
on every measurable set.
\begin{theorem}[Positive coefficient transfer]\label{thm:positive-transfer}
Let $l/T$ stay in a positive compact interval and let
$\boldsymbol B(z)=\sum_{d=0}^{D}z^d\nu_d$, where each $\nu_d$ is a
finite positive measure on any measurable space and $D\le C_0Lb$.
Then
\begin{equation}\label{eq:order-transfer}
 (1-\eps_n)r_{n,l}^{(b)}\boldsymbol B(1)
 \le [z^n]\{[u^l]R_b(z,u)\,\boldsymbol B(z)\}
 \le(1+\eps_n)r_{n,l}^{(b)}\boldsymbol B(1).
\end{equation}
In particular the bound holds after any collection of positive path
projections or weights. Its error does not depend on the dimension of
the space, the number of projections, or the mass of the event being
measured. If $\boldsymbol B(1)$ has positive total mass, normalizing
both sides to probability measures changes the multiplicative error only
by $1+O(\eps_n)$.
\end{theorem}
\begin{proof}
The middle measure is exactly
$\sum_{d=0}^D r_{n-d,l}^{(b)}\nu_d$.
Apply \eqref{eq:flatness} term by term. The last assertion follows by
applying the same inequalities to the whole space and dividing.
No norm bound for a product of projection operators is needed.

% BEGIN EXPANSION positive_normalization
\paragraph{Normalization and simultaneous restrictions.}
Here is the exact division used in the last assertion. Write
$\nu=\sum_d\nu_d$, $r=r_{n,l}^{(b)}$, and
$\mu=\sum_d r_{n-d,l}^{(b)}\nu_d$. For large $n$, $r>0$ and
$0\le\eps_n<1$. If $0<\nu(\Omega)<\infty$, the whole-space estimate
gives $0<\mu(\Omega)<\infty$. For every measurable $E$,
\[
 \frac{1-\eps_n}{1+\eps_n}\frac{\nu(E)}{\nu(\Omega)}
 \le\frac{\mu(E)}{\mu(\Omega)}
 \le\frac{1+\eps_n}{1-\eps_n}\frac{\nu(E)}{\nu(\Omega)}.
\]
If $\nu(E)=0$, both the middle and the reference event probabilities
are zero, so no division by an event probability is needed. More
generally a nonnegative measurable weight $w$ can be included by
replacing $\nu_d(E)$ with $\int_E w\dd\nu_d$, provided these weighted
measures are finite. A conjunction of path inequalities is exactly
such a weight, namely its indicator. All the inequalities are
imposed on the stored increment vector before forming this sum.

For the actual short-cycle application, fix the counts $q_i$ and let
$q=\sum_iq_i$, $l=k-q$. The product of the unnormalized harmonic
sample masses has total mass $\prod_i H_i^{q_i}/q_i!$. Grouping it
by its total size gives the measures $\nu_d$ above. Every possible
size is at most $qb\le kb$. Repeated samples and the factorial
$q_i!$ give precisely the cycle weights in
\eqref{eq:conditional-law}; on normalization this is the usual
multinomial expansion of the harmonic sampling law. Thus the
reference event $E$ can specify the entire multiblock path, its
endpoint, and restrictions on several observation points at once.
The one reservoir factor depends only on $d$ and $l$, irrespective
of how that event was described.
% END EXPANSION positive_normalization

\end{proof}

Here is its concrete use. Partition $\{1,\ldots,b\}$ into blocks
$I_0,\ldots,I_m$, and put $H_i=\sum_{j\in I_i}1/j$. For observation
points $t_1,\ldots,t_d$ let
$\boldsymbol f(j)=(\log|1-e(jt_v)|)_{v=1}^d$. A root is recorded as
$-\infty$ in the corresponding coordinate. Heights lie in
$(\R\cup\{-\infty\})^d$, with the usual absorbing convention for
addition; a positive exponential tilt assigns zero mass to such an atom.
After selecting counts $q_i$, a height-marked block is
\begin{equation}\label{eq:measure-block}
 \boldsymbol B_{i,q_i}(z)
 =\frac1{q_i!}\left(\sum_{j\in I_i}\frac{z^j}{j}
                       \delta_{\boldsymbol f(j)}\right)^{*q_i}.
\end{equation}
Convolution adds heights and size; before the final coefficient is taken,
one may retain the whole increment vector and impose all cumulative-height
inequalities. This gives positive measures with size degree at most
$(\sum q_i)b\le kb$. Evaluation at $z=1$, after normalization, is exactly
$q_i$ independent harmonic samples in each block, with mass $1/(jH_i)$.
Thus, conditionally on all $q_i$ and on $l=k-\sum q_i$ long cycles,
\eqref{eq:order-transfer} compares the \emph{entire} actual short-cycle
configuration with these product harmonic measures, with relative error
$1+O(\eps_n)$. The original total-size constraint has been integrated
out by its analytic coefficient, not replaced by a total-variation
approximation whose error is larger than the rare event.

\subsection{The exact-cycle distribution of the block counts}
\begin{proposition}[Multinomial count reservoir]\label{prop:multinomial}
Let $k/L=\kappa\in[\kappa_-,\kappa_+]$ and let
$(Q_0^*,\ldots,Q_m^*,L_*^*)$ be multinomial with $k$ trials and
probabilities
\begin{equation}\label{eq:multi-probs}
 H_0/L,\ldots,H_m/L,T/L.
\end{equation}
They sum to one. Under \eqref{eq:conditional-law}, let $Q_i$ count cycles
in $I_i$ and $L_*$ count cycles larger than $b$. Uniformly on
\begin{equation}\label{eq:l-typical}
 \mathcal G=\{|l-\kappa T|\le T^{3/4}\},
\end{equation}
the actual probability of each count vector equals the multinomial
probability times $1+o(1)$, with a uniform error. Both distributions assign
probability tending to one to $\mathcal G$. These assertions hold for
arbitrarily many blocks.
\end{proposition}
\begin{proof}
Sum the exact configuration weights at fixed $(q_i,l)$. All short
configurations have size at most $kb\le\kappa_+Lb$. By
Theorem~\ref{thm:positive-transfer}, their total weight is
\[
 \{1+O(\eps_n)\}\,r_{n,l}^{(b)}\prod_i\frac{H_i^{q_i}}{q_i!}.
\]
Divide by \eqref{eq:akn} and apply \eqref{eq:reservoir-marker}. The ratio
to the multinomial mass is
\[
 \frac{\Gamma(\kappa)}{\Gamma(l/T)}
       \{1+O(T^{-1}+\eps_n)\}=1+o(1)
\]
uniformly on \eqref{eq:l-typical}. A binomial Chernoff bound gives
$\Pp(|L_*^*-\kappa T|>T^{3/4})\le2e^{-c\sqrt T}$ uniformly.
Summing the uniform relative estimate on this set proves that the actual
mass there also tends to one. No assertion about the mass ratio outside
this set is necessary.
\end{proof}

\section{Pointwise killed-convolution estimates}
\label{sec:kernel}
All constants in this and the following sections can be chosen uniformly
for $\kappa\in[\kappa_-,\kappa_+]$. Thus $s=s_\kappa$ ranges over a
fixed compact subset of $(0,\infty)$.

\subsection{The Mellin kernel and a bridge density lower bound}
Let $V=\log|1-e(U)|$ for uniform $U\in\T$. A change of variables gives
\begin{equation}\label{eq:Vdensity}
 f_V(x)=\frac{e^x}{\pi\sqrt{1-e^{2x}/4}}\ind_{\{x<\log2\}}.
\end{equation}
For $\beta>0$, define the positive tilted kernel
\begin{equation}\label{eq:tilt-kernel}
 f_\beta(x)=e^{\beta x-\lambda(\beta)}f_V(x),\qquad
 \mu_\beta=\lambda'(\beta).
\end{equation}
Its convolution transform is $A(\beta+it)/A(\beta)$.

\begin{lemma}[Uniform convolution bounds]\label{lem:walk}
For $\beta$ in a fixed compact subset of $(0,\infty)$ and all sufficiently
large $q$, let $g_{\beta,q}$ be the continuous version of the centered
$q$-fold convolution density. Then
\[
 \sup_x g_{\beta,q}(x)\le Cq^{-1/2},\qquad
 g_{\beta,q}(0)\ge cq^{-1/2}.
\]
If $Y_j$ have density $f_\beta(\cdot+\mu_\beta)$ independently, then
\begin{equation}\label{eq:walkmax}
 \Pp\left(\max_{j\le q}\left|\sum_{v\le j}Y_v\right|>w\right)
 \le2e^{-cw^2/q},\qquad 1\le w\le c q.
\end{equation}
\end{lemma}
\begin{proof}
The centered log moment generating function is at most $Ct^2$ for
small real $t$, uniformly in $\beta$. The exponential supermartingale,
stopped at its first crossing before time $q$, gives
$\Pp(\max S_j>w)\le\exp(-tw+Cqt^2)$.
Choose $t$ proportional to $w/q$, and repeat with $-S_j$, proving
\eqref{eq:walkmax}.

For $q\ge3$, the convolution integral is continuous: $f_\beta$ belongs
to $L^{3/2}(\R)$, so Young's inequality and continuity of translations
give boundedness and continuity of its third and subsequent convolutions.
For the density bounds, the centered characteristic function has modulus
at most $e^{-ct^2}$ near zero, uniformly on the tilt interval. On each
fixed annulus its modulus is at most $\rho<1$: equality would force
$e^{itV}$ to be constant, impossible for the positive density
\eqref{eq:Vdensity} on an interval. The gamma ratio and Stirling's formula
on vertical strips give the bound $C(1+|t|)^{-1/2}$ at infinity.
Split the Fourier inversion integral into these three regions. The
absolute integral is at most $Cq^{-1/2}$. At zero, putting $t=u/\sqrt q$
in the first region gives
$\sqrt q\,g_{\beta,q}(0)\to(2\pi\lambda''(\beta))^{-1/2}$ uniformly;
the other regions are exponentially small. This proves both bounds.
\end{proof}

For a path of $q$ centered increments write $S_j$ for its partial sums.
The density of $S_q$ at zero restricted to
$\max_{j\le q}|S_j|\le W\sqrt q$ is well defined by convolution and
Lebesgue integration. All restricted densities below use the iterated
convolution integral with the relevant increment density, substituting
the prescribed total minus the preceding increments for the last increment.
\begin{lemma}[A pointwise bridge kernel]\label{lem:bridge}
There are fixed $W,c>0$ such that this restricted density is at least
$cq^{-1/2}$, uniformly for the same tilts and all sufficiently large $q$.
Consequently, if $d/q=\lambda'(\beta)$ and $\beta$ is in the same
compact tilt interval, the $f_s$ convolution density at total increment
$d$, restricted to
\[
 \left|\sum_{v\le j}x_v-\frac jq d\right|\le W\sqrt q
 \quad(1\le j\le q),
\]
is at least
\begin{equation}\label{eq:bridge-saddle}
 cq^{-1/2}\exp\{q[\lambda(\beta)-\lambda(s)]-(\beta-s)d\}.
\end{equation}
In particular, if $|d-q\lambda'(s)|/q=o(1)$, the last exponential is at
least $\exp\{-C(d-q\lambda'(s))^2/q\}$.
\end{lemma}
\begin{proof}
Split $q$ into $q_1=\lfloor q/2\rfloor$ and $q_2=q-q_1$.
The density at zero contributed by paths having
$\max_{j\le q_1}|S_j|>W\sqrt q$ is at most
\[
 \sup_x g_{\beta,q_2}(x)\,
 \Pp(\max_{j\le q_1}|S_j|>W\sqrt q)
 \le Cq^{-1/2}e^{-cW^2}.
\]
For a path with total zero, deviations during its second half are
negative reversed partial sums of that half. Reversing the increments
gives the same bound for a violation in the second half. Subtract the
two bounds from $g_{\beta,q}(0)\ge c_0q^{-1/2}$ and choose $W$ large.
These are statements about densities, obtained by integrating the unused
half-convolution; no conditioning on a null event is used without a density.

On the hyperplane of total increment $d$, the density ratio of the
$f_s$ product to the $f_\beta$ product is the constant exponential in
\eqref{eq:bridge-saddle}. The bridge tube is unchanged by that ratio.
Finally $\lambda''$ is bounded above and below on the compact tilt
interval. The inverse function theorem gives
$|\beta-s|\le C|d-q\lambda'(s)|/q$, and Taylor expansion of the
exponent gives the final assertion.

% BEGIN EXPANSION bridge_pointwise
\paragraph{The pointwise integral and its change of tilt.}
The definition used in this proof can be written explicitly. Given
increments $x_1,\ldots,x_{q-1}$, set
$x_q=d-\sum_{v<q}x_v$ and $S_j=\sum_{v\le j}x_v$. For a measurable
path restriction $E\subset\R^q$, put
\[
 K_{\beta,q}^{E}(d)=
 \int_{\R^{q-1}}\ind_E(x_1,\ldots,x_q)
              \prod_{v=1}^q f_\beta(x_v)
                   \dd x_1\cdots\dd x_{q-1}.
\]
This specifies a value at every $d$. With no restriction it agrees
with the continuous convolution density for $q\ge3$, by the
convolution identities and the regularity established in
Lemma~\ref{lem:walk}. Integrating in $d$ gives the probability of
$E$ under the product increment law: the substitution from the
last increment to $d$ is a translation with Jacobian one. The
same identity with an endpoint indicator gives disintegration
over any measurable endpoint set.

To apply the half-path argument, center the increments by
$\mu_\beta$. On the hyperplane of centered total zero, the
contribution of a first-half event $E_1$ is the integral of its
first-half product density times
$g_{\beta,q_2}(-S_{q_1})$. It is at most
$\sup g_{\beta,q_2}\,\Pp(E_1)$. For the second half use the reversal
$(x_1,\ldots,x_q)\mapsto(x_q,\ldots,x_1)$. On the free
$q-1$ coordinates its absolute Jacobian is one; its product density
is unchanged. Moreover a reversed partial sum of length $q-j$ is
$-S_j$ when the total is zero. A violation in either half is
therefore bounded by one of two identical maximal-inequality
estimates. Once $W$ is fixed, their sum is at most
$c_0/(2\sqrt q)$ by increasing $W$, and then increasing the lower
threshold for $q$ so that the allowed range in \eqref{eq:walkmax}
contains $W\sqrt q$. This gives a strictly positive bound at the
specified endpoint itself.

On the same parametrized hyperplane the change of tilt is exactly
\[
 \prod_{v=1}^q f_s(x_v)
 =\exp\{q[\lambda(\beta)-\lambda(s)]-(\beta-s)d\}
       \prod_{v=1}^q f_\beta(x_v).
\]
It remains true when one of the factors is zero. For the near-mean
assertion first enlarge the compact interval of $s$ slightly
within $(-1,\infty)$. On that larger interval let
$0<v_-\le\lambda''\le v_+$. For a fixed small enough $\delta>0$,
$|d/q-\lambda'(s)|\le\delta$ has a unique solution
$\lambda'(\beta)=d/q$ in the enlarged interval, and
$|\beta-s|\le v_-^{-1}|d/q-\lambda'(s)|$. Taylor's formula in the
form
\[
 \lambda(\beta)-\lambda(s)-(\beta-s)\lambda'(\beta)
 =-\int_s^\beta(t-s)\lambda''(t)\dd t
 \ge-\frac{v_+}{2}(\beta-s)^2
\]
now proves the asserted quadratic penalty, also when $\beta<s$.
The enlarged interval is essential when $s$ lies at one endpoint
of its original compact range.
% END EXPANSION bridge_pointwise

\end{proof}

This pointwise estimate is stronger than a lower bound for one endpoint
window. It will be integrated over a chain of endpoint boxes. The
integration is what prevents payment of a narrow-window factor at every
scale.

\subsection{Fine scales and a regular count environment}
Take a fixed $A_0$ sufficiently large and fixed constants $g_*,r_*>0$
to be chosen later. Define
\begin{equation}\label{eq:fine-scales}
 h=A_0\ell,\quad r=r_*\ell,\quad a_*=\log b,\quad
 m=\left\lceil\frac{a_*-r}{h}\right\rceil,\quad
 \omega=\frac{a_*-r}{m},\quad a_i=r+i\omega.
\end{equation}
Then $h/2\le\omega\le h$ for large $n$. Let $I_0$ contain $j\le e^r$
and $I_i$ contain $e^{a_{i-1}}<j\le e^{a_i}$, $1\le i\le m$,
with integer endpoints interpreted by these inequalities. They partition
$\{1,\ldots,b\}$. For the middle counts $q_i$ put
\[
 S_i=\sum_{j=1}^i q_j,\qquad Q=S_m.
\]
Their harmonic block masses satisfy
\begin{equation}\label{eq:harmonic-widths}
 H_i=\omega+O(e^{-a_{i-1}}),\qquad
 \sum_{i=1}^m|H_i-\omega|=O(e^{-r}).
\end{equation}

Fix $D_0>0$, to be chosen in Lemma~\ref{lem:boxes}, and put
\begin{equation}\label{eq:base-scale}
 H_*^{\rm c}=\frac{r^2}{D_0\log\ell},\qquad
 v_0=\left\lceil H_*^{\rm c}/\omega\right\rceil.
\end{equation}
Group the middle blocks dyadically from both ends. Start with $M=m$
remaining blocks and $v=v_0$. While $M\ge6v$, remove one group of $v$
blocks from each end of the remaining interval, and immediately replace
$(M,v)$ by $(M-2v,2v)$. At termination $2v\le M<6v$. Divide the
remaining blocks into $J_{\rm c}=\lceil M/(2v)\rceil$ consecutive groups,
each containing $\lfloor M/J_{\rm c}\rfloor$ or $\lceil M/J_{\rm c}\rceil$
blocks, placing the larger groups first. Order all groups from left to
right. Set $i_0=0$, and let $i_j$ be the last fine index of group $j$,
so that $i_B=m$. Write
\[
 \mathsf H_j=a_{i_j}-a_{i_{j-1}},\qquad
 \mathsf Q_j=S_{i_j}-S_{i_{j-1}},\qquad 1\le j\le B.
\]
For large $n$ this deterministic construction gives
\begin{equation}\label{eq:dyadic}
 c\ell\le B\le C\ell,\quad
 \mathsf H_j\ge H_*^{\rm c},\quad
 \mathsf H_1,\mathsf H_B\in[H_*^{\rm c},2H_*^{\rm c}],\quad
 1/5\le\mathsf H_{j+1}/\mathsf H_j\le5.
\end{equation}
The bounds follow by summing a geometric progression and noting that
$m/v_0\asymp L\log\ell/\ell^2$ with fixed-constant factors.
Define
\begin{equation}\label{eq:environment-E}
 E_j=1+\mathsf H_j^{-1/2}
 \max_{i_{j-1}\le i\le i_j}
 |S_i-S_{i_{j-1}}-\kappa(a_i-a_{i_{j-1}})|.
\end{equation}
For fixed $\eta>0,K_E,C_E$, call the environment regular if
\begin{equation}\label{eq:regular}
 \frac\kappa2\omega\le q_i\le(\kappa+\eta)\omega\quad(1\le i\le m),
 \qquad \max_jE_j\le K_E\sqrt{\log\ell},\qquad
 \sum_jE_j^2\le C_E B.
\end{equation}
All constants here can be selected uniformly over the compact $\kappa$
interval.

\begin{lemma}[Regularity under exact cycle conditioning]\label{lem:regular}
For each fixed $\eta>0$, constants $A_0,K_E,C_E$ can be chosen, independently
of fixed $r_*,D_0$, such that under $\Pp_{n,k}$, with probability tending
to one uniformly in $k$, the environment is regular,
\[
 Q_0\le C r,\qquad |L_*-\kappa T|\le T^{3/4}.
\]
In particular $L_*\le C\ell$. The threshold for $n$ may depend on the
subsequently fixed $r_*,D_0$.
\end{lemma}
\begin{proof}
By Proposition~\ref{prop:multinomial}, it suffices to prove these assertions
for the multinomial vector \eqref{eq:multi-probs}. Each $q_i$ is binomial
with mean $\kappa H_i$. A Chernoff bound and \eqref{eq:harmonic-widths},
followed by a union bound over at most $L$ blocks, establish the first
part of \eqref{eq:regular} when $A_0$ is large. The same bound gives
$Q_0\le Cr$ with probability tending to one. The assertion about $L_*$
was already proved.

Here are details for the multiscale energy; conditioning a fixed number
of trials must not be ignored. First form independent Poisson category
counts with the same means $k p_i$. On a coarse group, the centered
successive counts are independent-increment martingales. Exponential
martingales give
\[
 \Pp(E_j>x)\le C e^{-cx^2}
 \quad(2\le x\le c\sqrt{\mathsf H_j}),
\]
uniformly, with the negligible deterministic discrepancies in
\eqref{eq:harmonic-widths} absorbed. The union bound over $B=O(\ell)$
yields the desired maximum bound for a sufficiently large constant.
Doob's fourth-moment inequality and
$\E(\Poi(\mu)-\mu)^4=\mu+3\mu^2$ give
$\sup_j\E E_j^4\le C$. Different coarse groups are independent.
Hence Chebyshev's inequality shows $\sum E_j^2\le C_E B$ with failure
probability $O(B^{-1})$ when $C_E$ is sufficiently large.

To pass this elementary count estimate to a fixed number of trials,
take an independent $N\sim\Poi(k)$ and one sequence of iid category
labels of probabilities $(p_i)$. The first $N$ labels give the independent
Poisson counts, and the first $k$ give the multinomial counts. Conditional
on $D=|N-k|$, the differing labels are $D$ iid labels. If $D_j$ is their
number in coarse group $j$, then
\[
 \E\left[\sum_j\frac{D_j^2}{\mathsf H_j}\,\middle|\,D\right]
 \le C\left(\frac{DB}{L}+\frac{D^2}{L}\right).
\]
Indeed that category has probability
$\mathsf H_j/L+O(e^{-r}/L)$; use its binomial second moment and
\eqref{eq:dyadic}. Since $\E D\le\sqrt k$ and $\E D^2=k$, the last
sum is $O_{\Pp}(1)$. At every prefix of a coarse group the count change
is at most $D_j$. Thus the squared Euclidean distance between the two
vectors of $E_j$'s is $O_{\Pp}(1)$. Enlarging $K_E,C_E$ transfers their
maximum and energy bounds, since $\sqrt{\log\ell}\to\infty$ and
$B\to\infty$. This use of Poisson variables is only a proof of a
multinomial count inequality. There is no comparison of spectral fields
or rare high-point probabilities at this step.

% BEGIN EXPANSION count_coupling
\paragraph{Quantitative effect of fixing the number of trials.}
The preceding coupling controls all prefixes simultaneously.
Write $p_j^{\rm c}$ for the probability that one category label
lies in coarse group $j$. For large $n$,
$p_j^{\rm c}\le C\mathsf H_j/L$ and
$\sum_jp_j^{\rm c}\le1$. Conditional on $D$, each $D_j$ is
binomial with parameters $D,p_j^{\rm c}$, even though the $D_j$
are not independent. Hence, without using independence among them,
\[
 \begin{split}
 \E\left[\sum_j\frac{D_j^2}{\mathsf H_j}\,\middle|\,D\right]
 &\le D\sum_j\frac{p_j^{\rm c}}{\mathsf H_j}
       +D^2\sum_j\frac{(p_j^{\rm c})^2}{\mathsf H_j}\\
 &\le CDB/L+CD^2/L.
 \end{split}
\]
Since $N$ is independent of the entire label sequence, the labels
between indices $\min(N,k)+1$ and $\max(N,k)$ remain iid after
conditioning on $N$, hence after conditioning on $D$ for this
calculation. Also $\E D^2=\Var N=k$ and
$\E D\le\sqrt{k}$. The expectation of this energy is bounded
uniformly because $B=O(\ell)$ and $k\asymp L$.

Let $\widetilde E_j$ and $E_j$ be the Poisson and multinomial
environments in this one coupling, with the same deterministic
centering. The elementary inequality
$|\max_i|u_i|-\max_i|v_i||\le\max_i|u_i-v_i|$ gives
\[
 |E_j-\widetilde E_j|\le D_j/\sqrt{\mathsf H_j},
 \qquad
 \sum_j E_j^2\le2\sum_j\widetilde E_j^2
                         +2\sum_jD_j^2/\mathsf H_j.
\]
Markov's inequality makes the second sum $o_{\Pp}(B)$, and its
square root $o_{\Pp}(\sqrt{\log\ell})$. Enlarging the two fixed
constants in the Poisson bounds therefore yields both parts of
\eqref{eq:regular} for the multinomial vector. No conditioning
estimate is divided by $\Pp(N=k)$, which would lose a factor
of order $\sqrt L$.
% END EXPANSION count_coupling

\end{proof}

\subsection{Dyadic boxes instead of an assumed barrier theorem}
Let $G=g_*\ell$, $s=s_\kappa$, and $\lambda=\lambda(s)$. Define
\begin{equation}\label{eq:frontier}
 F_i=\frac{a_i+\lambda S_i}{s},\qquad
 y=\frac{a_*-r+\lambda Q}{s},\qquad B_i=F_i-G.
\end{equation}
In each fine block convolve $q_i$ copies of $f_s$ and one independent
noise of modulus at most $L^{-10}$. Write $W_i$ for the cumulative height
through block $i$. Let
\begin{equation}\label{eq:path-weight}
 \mathcal A=\{W_i\le B_i\ (1\le i\le m),\quad y\le W_m\le y+1\},
 \qquad
 p_L=\E_s[e^{-s(W_m-y)}\ind_{\mathcal A}].
\end{equation}
This expectation is simply the integral of a product of the positive
convolution kernels, with the stated projections.

\begin{lemma}[Polynomial lower mass for the constrained kernel]
\label{lem:boxes}
There exist $D_0,C_*<\infty$, depending on the compact $\kappa$ interval
and on $A_0,K_E,C_E$, with the following property. For all fixed $g_*,r_*$
satisfying
\begin{equation}\label{eq:rg}
 r_*\ge64 s_{\max}g_*+1,\qquad
 s_{\max}=\max_{\kappa\in[\kappa_-,\kappa_+]}s_\kappa,
\end{equation}
and all regular environments, uniformly for large $n$,
\begin{equation}\label{eq:polynomial-cost}
 p_L\ge L^{-C_*}.
\end{equation}
The constants $D_0,C_*$ do not depend on the fixed $g_*,r_*$. The
threshold for $n$ is allowed to depend on them.
\end{lemma}
\begin{proof}
We exhibit a lower bound for a product of killed kernels; no adaptive
random-environment measure is an input. Temporarily omit the noises.
Choose a large constant $B_0$. At a nonterminal coarse endpoint define
\[
 \gamma_j=4G+B_0(\sqrt{\mathsf H_j}E_j+
                         \sqrt{\mathsf H_{j+1}}E_{j+1}),
 \qquad z_j=F_{i_j}-\gamma_j \quad(1\le j<B).
\]
At the two ends take
\[
 z_0=0,\quad\gamma_0=r/s;\qquad
 z_B=y+\tfrac12,\quad\gamma_B=r/s-\tfrac12.
\]
Let the intermediate endpoint boxes be
$\mathcal I_j=[z_j-\sqrt{\mathsf H_j},z_j+\sqrt{\mathsf H_j}]$,
$1\le j<B$, and let $\mathcal I_B=[z_B-1/10,z_B+1/10]$.
The initial box is the singleton $\{0\}$.

Choose $B_0$ to dominate the bridge width in Lemma~\ref{lem:bridge},
the neighboring-width ratio five, and the constants in the next
interpolation calculation. Choose $D_0$ so large that
\[
 C B_0 K_E\sqrt{2/D_0}<1/(8s_{\max}),
\]
where $C$ is a fixed constant large enough for those same bounds.
By \eqref{eq:rg}, both endpoint gaps then dominate $4G$ plus the required
multiple of $E_1\sqrt{\mathsf H_1}$ or
$E_B\sqrt{\mathsf H_B}$, for large $n$. Indeed these products are at
most $K_E r\sqrt{2/D_0}$, whereas $r/s\ge r/s_{\max}$.

Consider group $j$ and any $x\in\mathcal I_{j-1}$,
$x'\in\mathcal I_j$. Put $d=x'-x$ and $q=\mathsf Q_j$.
The critical identity $s\lambda'(s)-\lambda=1/\kappa$ gives
\begin{equation}\label{eq:box-drift}
 d-q\lambda'(s)
 =\frac{\mathsf H_j-q/\kappa}{s}
   -(\gamma_j-\gamma_{j-1})+(e_j-e_{j-1}),
\end{equation}
where $e_j=x'-z_j$ and $e_{j-1}=x-z_{j-1}$.
Use $E_0=E_{B+1}=1$ when needed. Uniformly over these boxes,
\begin{equation}\label{eq:Ub}
 |d-q\lambda'(s)|\le C\sqrt{\mathsf H_j}\,U_j,
 \quad
 U_j=1+E_{j-1}+E_j+E_{j+1}
       +\frac r{\sqrt{\mathsf H_j}}\ind_{\{j=1\text{ or }j=B\}}.
\end{equation}
For interior groups the $4G$ terms cancel exactly. At the boundary they
are bounded by a constant times $r$. Moreover $q\asymp\mathsf H_j$,
$\max_j U_j/\sqrt{\mathsf H_j}=O(\log\ell/r)=o(1)$, and
\begin{equation}\label{eq:box-energy}
 \sum_{j=1}^B U_j^2
 \le C\left(B+\sum_jE_j^2+\frac{r^2}{H_*^{\rm c}}\right)
 \le C\ell.
\end{equation}
The last constant is independent of $g_*,r_*$, since
$r^2/H_*^{\rm c}=D_0\log\ell$.

We verify the path restriction before using its density. At a fine
endpoint $i$ in this group, set
$f=(S_i-S_{i_{j-1}})/q$. The straight line in sample count joining $x$
to $x'$ has gap below $F_i$ equal to
\begin{equation}\label{eq:chord-gap}
 (1-f)(\gamma_{j-1}-e_{j-1})+f(\gamma_j-e_j)
       +\frac{a_i-a_{i_{j-1}}-f\mathsf H_j}{s}.
\end{equation}
If $D_i=S_i-S_{i_{j-1}}-\kappa(a_i-a_{i_{j-1}})$, the numerator
in its final term is
\[
 \frac{(a_i-a_{i_{j-1}})D_{i_j}-\mathsf H_jD_i}{q}.
\]
It has modulus at most $CE_j\sqrt{\mathsf H_j}$.
Our choices of gaps, $B_0,D_0$ make \eqref{eq:chord-gap} exceed
$2G+W\sqrt q$ at every such endpoint. Hence a bridge lying within
$W\sqrt q$ of that straight line stays below $F_i-2G$ at all fine
endpoints.

For each fixed pair $x,x'$, solve the deterministic saddle equation
$\lambda'(\beta)=d/q$. It has a solution near $s$, uniformly by
\eqref{eq:Ub}, and the compactness needed in Lemma~\ref{lem:bridge}
holds. That lemma gives the following \emph{pointwise} lower bound for
the group convolution density, killed at any violated fine barrier:
\begin{equation}\label{eq:killed-minorization}
 K_j(x,x')\ge c\mathsf H_j^{-1/2}\exp(-CU_j^2),
 \qquad x\in\mathcal I_{j-1},\quad x'\in\mathcal I_j.
\end{equation}
This step chooses a saddle for each density evaluation; it does not
assert that the random increments follow a preassigned adaptive law.

Multiply \eqref{eq:killed-minorization} and integrate all intermediate
endpoints over their boxes. Their widths $2\sqrt{\mathsf H_j}$ cancel
the respective density factors. Only the last, fixed-width endpoint
window leaves a factor $\mathsf H_B^{-1/2}$. Thus the mass of paths
satisfying the stronger barriers and ending in $\mathcal I_B$ is at least
\[
 c^B\mathsf H_B^{-1/2}\exp\left(-C\sum_jU_j^2\right)\ge L^{-C_*}.
\]
Indeed $B=O(\ell)$, \eqref{eq:box-energy} holds, and
$\log\mathsf H_B=O(\log\ell)+O(\log r_*)=o(\ell)$ for each fixed $r_*$.
Finally all noises together change any height by at most $mL^{-10}\le
L^{-9}$. The strict barrier margins and the endpoint interval
$[y+2/5,y+3/5]$ absorb them. On \eqref{eq:path-weight} the terminal
weight is at least $e^{-s}$. Increasing $C_*$ proves
\eqref{eq:polynomial-cost}.

% BEGIN EXPANSION box_integration
\paragraph{Endpoint costs and the order of constants.}
We give the bookkeeping behind the two cancellations just used.
From regularity,
$|q-\kappa\mathsf H_j|\le E_j\sqrt{\mathsf H_j}$.
Furthermore
$\max_j E_j/\sqrt{\mathsf H_j}
 \le K_E\sqrt{D_0}\log\ell/r=o(1)$.
Thus $q/\mathsf H_j$ lies, for all large $n$, in a positive compact
interval determined by the original $\kappa$ interval. This bound
comes from the coarse environment energy; it does not require
any further restriction on the fine-count upper tolerance.

The interior gap differences in \eqref{eq:box-drift} involve at
most three neighboring $E$'s. Adjacent square-root widths have
ratios between $1/\sqrt5$ and $\sqrt5$. The inequality
$(t_1+\cdots+t_5)^2\le5\sum_{\nu=1}^5t_\nu^2$, applied to
\eqref{eq:Ub} and then summed in $j$, consequently gives
\[
 \sum_jU_j^2
 \le C\left(B+\sum_jE_j^2+
                     r^2/\mathsf H_1+r^2/\mathsf H_B\right)
 \le C(1+C_E)B+2CD_0\log\ell.
\]
Every interior $E_j^2$ is counted at most three times. Also
$U_j/\sqrt{\mathsf H_j}=o(1)$ uniformly, including the two
boundary terms $r/\mathsf H_j\le D_0\log\ell/r$. These bounds
justify use of the fixed near-mean neighborhood from
Lemma~\ref{lem:bridge} for every pair of box endpoints.

For the path inclusion, write
$\tau=a_i-a_{i_{j-1}}$ and $D_i=S_i-S_{i_{j-1}}-\kappa\tau$.
Since $q=\kappa\mathsf H_j+D_{i_j}$, elementary cancellation gives
\[
 \tau-\frac{\kappa\tau+D_i}{q}\mathsf H_j
   =\frac{\tau D_{i_j}-\mathsf H_jD_i}{q}.
\]
Here $0\le\tau\le\mathsf H_j$ and both deviations are at most
$E_j\sqrt{\mathsf H_j}$ in modulus. This is the error charged
against the two adjacent endpoint gaps in \eqref{eq:chord-gap}.
At the initial and final endpoints the additional gap $r/s$
supplies the same margin by the choice of $D_0$ and \eqref{eq:rg}.
The final endpoint is centered at $y+1/2$, so its gap is
$r/s-1/2$, exactly as specified above.

Tonelli's theorem and the independence of disjoint groups identify
the constrained endpoint-chain mass with the iterated integral
of their killed kernels. The quantitative lower bound for that
integral is
\[
 \left(\prod_{j=1}^B c\mathsf H_j^{-1/2}e^{-CU_j^2}\right)
 \left(\prod_{j=1}^{B-1}2\sqrt{\mathsf H_j}\right)\frac15
 =\frac{2^{B-1}c^B}{5}\,
       \mathsf H_B^{-1/2}e^{-C\sum_jU_j^2}.
\]
Taking logarithms leaves a loss at most
$C'B+C''\sum_jU_j^2+\tfrac12\log\mathsf H_B+C'''$.
The coefficient of $\ell$ in this bound can be fixed before
$g_*,r_*$: $B\le C\ell$ with a fixed coefficient, the energy is
bounded as above, and
$\log\mathsf H_B=O(\log\ell)+O(|\log r_*|+|\log D_0|)$.
Only the threshold for $n$ changes when these fixed parameters
change. Finally one can first fix any noise vector satisfying
the coordinate bound, apply the deterministic margin argument,
and then integrate that vector. Thus this lower bound even
allows dependence among the noise coordinates; the noise vector
is still sampled independently of the increment path.
% END EXPANSION box_integration

\end{proof}

\begin{remark}[Why the exponent is logarithmic]
A separate fixed-width endpoint condition in every coarse group would
produce $\prod_j\mathsf H_j^{-1/2}$ and a much larger logarithmic cost.
The pointwise kernel estimate permits boxes of width
$\sqrt{\mathsf H_j}$ and cancels all intermediate factors by integration.
Only $\sum E_j^2$, not $B\max E_j^2$, is charged. These two facts are
essential to the precision in Theorem~\ref{thm:main}.
\end{remark}

\section{Fourier remainders for the marked height kernels}
\label{sec:fourier}
These estimates concern only the explicitly evaluated kernels
\eqref{eq:measure-block}. The coefficient-reservoir theorem is what
allows them to be used under the original size and cycle constraints.

\subsection{Fourier regularity and arithmetic separation}
For $\Ree z>0$ put $\phi_z(t)=|1-e(t)|^z$, set its value at an integer
to zero, and use the real logarithm of $|1-e(t)|$ off the integers.
Use the convention $\widehat\phi_z(j)=\int_\T\phi_z(t)e(-jt)\dd t$.
For $s>0$ its Fourier coefficients are
\begin{equation}\label{eq:fourier-real}
 \widehat\phi_s(j)
 =\frac{(-1)^j\Gamma(1+s)}
 {\Gamma(1+s/2-j)\Gamma(1+s/2+j)}.
\end{equation}
The formula follows from the beta integral; alternatively the
denominator-free recurrence
$(s/2+j)\widehat\phi_s(j)=-(1+s/2-j)\widehat\phi_s(j-1)$ for $j\ge1$,
together with $\widehat\phi_s(0)=A(s)$ and
$\widehat\phi_s(-j)=\widehat\phi_s(j)$, determines the coefficients.
Reciprocal gamma zeros are interpreted literally.

\begin{lemma}[Uniform Fourier tail]\label{lem:fourier}
Fix $0<p_-<p_+<\infty$ and $0<\alpha<\min(p_-,1)$. For
$p\in[p_-,p_+]$, $u\in\R$, and every $R>0$,
\begin{equation}\label{eq:fourier-tail}
 \sum_j|\widehat\phi_{p+iu}(j)|\le C(1+|u|)^2,
 \qquad
 \sum_{|j|>R}|\widehat\phi_{p+iu}(j)|
 \le C(1+|u|)^2 R^{-\alpha}.
\end{equation}
\end{lemma}
\begin{proof}
The function is absolutely continuous and its derivative is integrable,
since near an integer its magnitude is bounded by
$C(1+|u|)|t|^{p_--1}$. Split an $L^1$ translation difference of the
derivative at distance $h$ from the integers. Within that distance its
integral is at most $C(1+|u|)h^{p_-}$. Away from it, integrate the second
derivative bound $C(1+|u|)^2(|t|^{p_--2}+1)$ over the translation path.
This gives
\[
 \|\phi'_{p+iu}(\cdot+h)-\phi'_{p+iu}\|_1
 \le C(1+|u|)^2 h^\alpha.
\]
The smaller exponent $\alpha$ absorbs the logarithm at $p_-=1$.
Use $h=1/(2|j|)$ in the Fourier coefficient of this difference, and
$\widehat{\phi'}(j)=2\pi ij\widehat\phi(j)$. It follows that
$|\widehat\phi_{p+iu}(j)|\le C(1+|u|)^2|j|^{-1-\alpha}$ for $j\ne0$.
The zero coefficient is bounded by $A(p)$. Summation proves the lemma.
\end{proof}

For a nonempty harmonic block $I(a,c)=\{j:e^a<j\le e^c\}$ let
$H=\sum_{j\in I(a,c)}1/j$ and define
\[
 B_{a,c}(\boldsymbol z;\boldsymbol t)
 =\frac1H\sum_{j\in I(a,c)}\frac1j
                   \prod_{v=1}^d\phi_{z_v}(jt_v),\qquad d=1,2.
\]
Summation by parts in a geometric progression yields
\begin{equation}\label{eq:abel}
 \left|\sum_{j\in I(a,c)}\frac{e(ju)}j\right|
       \le\frac{Ce^{-a}}{\|u\|_\T}\quad(u\notin\Z).
\end{equation}
It follows from Lemma~\ref{lem:fourier} that whenever
\begin{equation}\label{eq:separation}
 \|k_1t_1+\cdots+k_dt_d\|_\T\ge e^{-a+\Delta}
 \quad(0\ne\boldsymbol k\in\Z^d,\ |k_v|\le R),
\end{equation}
$\Ree z_v\in[p_-,p_+]$, and $|\Im z_v|\le T_0$, one has
\begin{equation}\label{eq:B-error}
 \left|B_{a,c}(\boldsymbol z;\boldsymbol t)-\prod_v A(z_v)\right|
 \le C(1+T_0)^{2d}\left(R^{-\alpha}+\frac{e^{-\Delta}}H\right).
\end{equation}
To see this, expand the absolutely convergent Fourier series. The zero
vector contributes the stated product. On the box $|k_v|\le R$ use
\eqref{eq:abel}; off it use one tail estimate and the remaining absolute
Fourier norms. This also proves uniformity when the observation points
move with $n$, subject only to \eqref{eq:separation}.

\subsection{Smoothed kernel comparison with relative-event accuracy}
\begin{lemma}[Polynomially accurate kernel comparison]\label{lem:smooth}
For each fixed $J>0$ there are positive constants $u_1,u_2,u_3$ such that
the following holds, uniformly in the compact $\kappa$ interval. Put
\begin{equation}\label{eq:smooth-params}
 \delta=L^{-10},\quad T_0=L^{u_1},\quad
 R=\lceil L^{u_2}\rceil,\quad\Delta=u_3\ell.
\end{equation}
In a fine block with $q\asymp\omega\asymp\ell$, take $q$ harmonic samples
and their $d$ log-sine sums $X_v$, $d=1,2$. Suppose \eqref{eq:separation}
holds. Tilt their joint height measure by $e^{s\sum_vX_v}$ and normalize.
Add independently to each coordinate a noise $\zeta$ which is the sum
of ten uniforms on $[-\delta/10,\delta/10]$. This normalized smoothed
measure differs in total variation by at most $L^{-J}$ from $d$
independent copies of $f_s^{*q}$ convolved with that noise. Also
\begin{equation}\label{eq:normalizer}
 B_{a,c}(s,\ldots,s;\boldsymbol t)^q
       =A(s)^{dq}\{1+O(L^{-J})\}.
\end{equation}
\end{lemma}
\begin{proof}
Choose a positive $\alpha<\min(s_{\min}/2,1)$, and apply
\eqref{eq:B-error} to a real tilt interval containing $s/2,s$ and $3s/2$
for every possible $s$. The characteristic function of the tilted vector
is
\[
 \left\{\frac{B_{a,c}(s+it_1',\ldots,s+it_d';\boldsymbol t)}
 {B_{a,c}(s,\ldots,s;\boldsymbol t)}\right\}^{q}.
\]
On $[-T_0,T_0]^d$ compare it with
$\{\prod_v A(s+it_v')/A(s)^d\}^{q}$.
The normalized quantities have modulus at most one, so their $q$th
powers differ by at most
$Cq(1+T_0)^{2d}(R^{-\alpha}+e^{-\Delta}/H)$.
The real denominator is bounded away from zero for large $n$.

The noise transform has modulus at most
$\min(1,C(\delta|t|)^{-10})$. Fourier inversion bounds the uniform
density error inside the frequency box by
$CT_0^d q(1+T_0)^{2d}(R^{-\alpha}+e^{-\Delta}/H)$; outside it the
bound is $C\delta^{-d}(\delta T_0)^{-9}$. For example, first choose
$u_1$ with $9u_1>J+160$, then $u_2,u_3$ with
\[
 \alpha u_2>6u_1+J+30,\qquad u_3>6u_1+J+30.
\]
The density error is then $O(L^{-J-10})$.
On $[-L^2,L^2]^d$ integration costs at most $O(L^4)$. Outside this
spatial box positive tails vanish, since $X_v\le q\log2$. Under the
tilt, a negative tail is at most
\[
 e^{-sL^2/3}
 \left\{\frac{B_{a,c}(s/2,s,\ldots,s;\boldsymbol t)}
 {B_{a,c}(s,\ldots,s;\boldsymbol t)}\right\}^{q}
 \le e^{-cL^2}.
\]
The ratio is bounded by \eqref{eq:B-error}; the reference measures have
the same tail bound. This proves the total variation assertion.
\eqref{eq:normalizer} follows from the real version of \eqref{eq:B-error},
increasing the fixed exponents if necessary. Atoms at $-\infty$ have
zero tilted mass. No total variation comparison with a continuous kernel
is asserted before smoothing.

% BEGIN EXPANSION smoothing_budget
\paragraph{Frequency tails, normalization, and product errors.}
To make the frequency estimate explicit, let $h_\delta$ be the
characteristic function of the sum of the ten uniforms. Direct
integration of a uniform variable gives
\[
 h_\delta(t)=
 \left(\frac{\sin(\delta t/10)}{\delta t/10}\right)^{10},
 \qquad
 \int_\R|h_\delta(t)|\dd t\le C\delta^{-1}.
\]
The ratio at zero is interpreted as one. For $\delta T_0\ge1$,
\[
 \int_{|t|>T_0}|h_\delta(t)|\dd t
 \le C\delta^{-1}(\delta T_0)^{-9}.
\]
For $d=1,2$, the complement of the frequency cube is covered by
the $d$ events $|t_v|>T_0$. Integrating the product noise transform
and using these two bounds gives
$C\delta^{-d}(\delta T_0)^{-9}$. Both normalized characteristic
functions have modulus at most one, which justifies this bound
for their difference as well. With $\delta=L^{-10}$ its worst
exponent, at $d=2$, is $110-9u_1$. Inside the cube the exponent is
at most $6u_1+1-\min(\alpha u_2,u_3)$, using $q\le L$ and $H\ge1$.
The stated choices make both bounds smaller than $L^{-J-10}$
after absorbing fixed constants. The integration volume
$O(L^4)$ then remains within the total variation budget.

The normalizer is controlled separately from this total
variation estimate. If a real block normalizer is
$A(s)^d(1+e_i)$, its relative error tends to zero uniformly, and
$|e_i|\le1/2$ for large $n$. The inequality
$|\log(1+e_i)|\le2|e_i|$ shows that raising it to the $q_i$th
power costs at most $2q_i|e_i|$ in its logarithm. The frequency
exponents can be increased once so that the desired
$L^{-J}$ bound includes this factor. On multiplying the fine
block normalizers the logarithmic errors add.

For normalized probability kernels $\mu_i,\nu_i$, changing
one factor at a time in their product proves
\[
 \left|\int F\dd\bigotimes_{i=1}^m\mu_i
       -\int F\dd\bigotimes_{i=1}^m\nu_i\right|
 \le\sum_{i=1}^m\|\mu_i-\nu_i\|_{\TV},
 \qquad 0\le F\le1.
\]
One may take total variation here to be the supremum over
tests in $[0,1]$; the $L^1$ density bound also bounds this
quantity. In each telescoping term the other factors are
probability measures, so the same inequality applies however
many coordinates the path test $F$ examines. For the present
kernels the resulting bound is $mL^{-J}$, with the same
choice of exponents for one and two observation points.
% END EXPANSION smoothing_budget

\end{proof}

\subsection{An integrated bound for every arithmetic arc}
For $s>0$ define
\begin{equation}\label{eq:bs}
 b_s=\begin{cases}0,&0<s\le2,\\
 (s-1)\log2-\lambda(s),&s>2.
 \end{cases}
\end{equation}
\begin{lemma}[Entropy gap and a block envelope]\label{lem:entropy}
At $s=s_\kappa$ one has $\kappa b_s<1$, with a uniform strict gap on
the compact $\kappa$ interval. For every block,
\begin{equation}\label{eq:envelope}
 B_{a,c}(s;t)\le A(s)\exp\{b_s+C/H\}\quad(t\in\T).
\end{equation}
\end{lemma}
\begin{proof}
For an arbitrary interval of positive integers,
$\sum\cos(2\pi ju)/j\ge-C$, uniformly in $u$. Indeed the cosines are
nonnegative up to $(6\|u\|_\T)^{-1}$; the remaining sum is bounded by
\eqref{eq:abel} starting at that cutoff. If the cutoff is below one,
start at one, and if $u$ is integral the sum is positive.
For $0<s\le2$, all nonzero coefficients in \eqref{eq:fourier-real} are
nonpositive. Pairing them and applying the preceding lower bound gives
$B_{a,c}(s;t)\le A(s)+C/H$. For $s>2$ use
$\phi_s\le2^{s-2}\phi_2$, giving $B_{a,c}(s;t)\le2^{s-1}+C/H$.
This proves \eqref{eq:envelope}.

The gamma duplication formula and the integral for the difference of two
digamma functions give
\[
 \log2-\lambda'(s)=\int_0^\infty\frac{e^{-sx}}{1+e^x}\dd x,
 \qquad 0<s(\log2-\lambda'(s))<\tfrac12.
\]
For example, the integral follows by subtracting the two logarithmic
derivatives in the duplication formula and expanding
$(1+e^x)^{-1}$, or directly integrating the convergent digamma difference.
Thus for $s>2$, \eqref{eq:critical} gives
\[
 \kappa b_s=1+\kappa\{s(\log2-\lambda'(s))-\log2\}<1.
\]
For $s\le2$ the assertion is immediate. Continuity and compactness make
the gap uniform.
\end{proof}

Choose $\eta>0$ so small that $(\kappa+\eta)b_{s_\kappa}<1$ uniformly.
Use this value in \eqref{eq:regular}. If $X_i(t)$ denotes the harmonic
height sum in fine block $i$ and $X_{\rm mid}=\sum_iX_i$, then
\begin{lemma}[Integrated middle moment]\label{lem:integrated}
For sufficiently large fixed $r_*$ (as specified below), uniformly over
regular count vectors,
\begin{equation}\label{eq:integrated}
 \E\int_\T e^{sX_{\rm mid}(t)}\dd t\le A(s)^Q\{1+o(1)\}.
\end{equation}
The expectation here is evaluation of the normalized product of harmonic
kernels, not an assumption about independence in the original Ewens law.
\end{lemma}
\begin{proof}
Let $d_R(t)=\min_{1\le|j|\le R}\|jt\|_\T$.
For $d_R(t)\ge e^{-r+\Delta}$ every block is separated and
\eqref{eq:normalizer} gives $A(s)^Q(1+o(1))$.
Partition the other points according to
\[
 e^{-a_k+\Delta}\le d_R(t)<e^{-a_{k-1}+\Delta},\quad1\le k\le m,
\]
including all smaller distances in class $m$. That class has measure at
most $CR e^{-a_{k-1}+\Delta}$. Blocks after $k$ are separated.
Using \eqref{eq:envelope}, $q_i\le(\kappa+\eta)\omega$, and
$\omega\to\infty$, there is a fixed $\rho<1$ such that
\[
 \prod_{i\le k}B_{a_{i-1},a_i}(s;t)^{q_i}
 \le A(s)^{S_k}e^{\rho(a_k-r)}.
\]
The per-block $Cq_i/H_i$ errors contribute $O(k)$, which is absorbed
by the uniform strict gap since $a_k-r=k\omega$.
The suffix contributes $A(s)^{Q-S_k}(1+o(1))$; no suffix is required for
$k=m$. The normalized integral over all these classes is at most
\[
 C L^{u_2+u_3}e^{-r+\rho\omega}
 \sum_{k=1}^m e^{-(1-\rho)(a_{k-1}-r)}=o(1)
\]
when $r_*>u_2+u_3+A_0+2$. The geometric sum is bounded.
\end{proof}

\section{Closing the maximum bounds inside the exact coefficients}
\label{sec:maximum}

\subsection{Order of constants}
All choices are uniform over the fixed interval of $\kappa$. Choose
$\eta$ from Lemma~\ref{lem:entropy}, then $A_0,K_E,C_E$ from
Lemma~\ref{lem:regular}. Select $D_0,C_*$ in Lemma~\ref{lem:boxes}.
Next fix $J>2C_*+12$ and select $u_1,u_2,u_3$ from
Lemma~\ref{lem:smooth}. Finally take $g_*,r_*$ so large that
\begin{align}\label{eq:constant-order}
 s_{\min}g_*&>2u_2+u_3+A_0+2C_*+12,\\
 r_*&>64s_{\max}g_*+1,\qquad
 r_*>2u_2+u_3+A_0+12.\nonumber
\end{align}
This is not circular: Lemma~\ref{lem:boxes} explicitly made $C_*$
independent of the last two constants. Throughout this section the count
vector is regular, until its probability is restored at the end.

\subsection{A first and relative second kernel calculation}
Let the harmonic product measure at fixed counts be denoted by $\E_H$.
Define
\[
 \mathcal D_L=\{t\in\T:\|jt\|_\T\ge e^{-r+\Delta},\ 1\le |j|\le R\}.
\]
Its complement has measure at most $CR e^{-r+\Delta}=o(1)$.
The following lower bound is uniform over every measurable
$\mathcal D\subset\mathcal D_L$ of measure at least $1/2$.
Let $\zeta_i$ be independent noises as in Lemma~\ref{lem:smooth}, put
$U_\zeta=\sum_i\zeta_i$, and define a nonnegative functional of the
actual harmonic samples by
\begin{equation}\label{eq:Z}
 Z(t)=\E_\zeta\left[e^{-sU_\zeta}
  \ind_{\{(\sum_{j\le i}(X_j(t)+\zeta_j))_{i\le m}\in\mathcal A\}}
                \right],\qquad Z_{\mathcal D}=\int_{\mathcal D}Z(t)\dd t.
\end{equation}
The event $\mathcal A$ is exactly \eqref{eq:path-weight}; the notation
means its cumulative-height inequalities. If $Z_{\mathcal D}>0$, then
some actual point in $\mathcal D$ has middle height at least $y-m\delta$.
The noises are auxiliary integrations, not alterations of the final
spectral field.

For $t\in\mathcal D_L$, tilt each fine height kernel by $e^{sX_i(t)}$.
The normalizer is $A(s)^Q(1+o(1))$. The noise factor in \eqref{eq:Z}
combines with the inverse tilt to give the terminal weight
$e^{-s\sum_i(X_i(t)+\zeta_i)}$. After removing $e^{-sy}$, the test
function is bounded by one on $\mathcal A$.
Telescoping the product of the normalized smoothed measures gives error
at most $mL^{-J}$. Since $m\le L$ and $p_L\ge L^{-C_*}$,
\begin{equation}\label{eq:rare-accuracy}
 mL^{-J}=o(p_L^2),
\end{equation}
not just $o(1)$. Lemma~\ref{lem:smooth} therefore gives uniformly in $t$
\begin{equation}\label{eq:first-moment}
 \E_H Z(t)=(1+o(1))\mu_Lp_L,\qquad
 \mu_L=e^{-sy+\lambda Q}=e^{-a_*+r}.
\end{equation}
Consequently $\E_H Z_{\mathcal D}=(1+o(1))\Leb(\mathcal D)\mu_Lp_L$.

For a pair define
\[
 d_R(t,u)=\min_{\substack{(j,l)\in\Z^2\setminus\{(0,0)\}\\|j|,|l|\le R}}
                         \|jt+lu\|_\T.
\]
If $d_R(t,u)\ge e^{-r+\Delta}$, apply the two-point comparison in every
block. The reference kernels at the two points are independent, including
independent noise vectors in the two copies of \eqref{eq:Z}. Thus
\begin{equation}\label{eq:far-pair}
 \E_H[Z(t)Z(u)]=(1+o(1))\mu_L^2p_L^2
\end{equation}
uniformly over these far pairs. Equation~\eqref{eq:rare-accuracy} is
what makes this relative statement valid.

% BEGIN EXPANSION first_far_moments
\paragraph{The bounded tests after the change of measure.}
Let $\mathcal N_t$ be the product of the one-point block
normalizers and let $\widetilde\Pp_t$ be the corresponding
product tilted law, before adjoining the noises. Set
$\widetilde W_i=\sum_{j\le i}(X_j(t)+\zeta_j)$. The exact
identity, with no asymptotic comparison yet, is
\[
 \E_H Z(t)=
 e^{-sy}\mathcal N_t\,
 \widetilde\E_{t,\zeta}
       [e^{-s(\widetilde W_m-y)}\ind_{\mathcal A}].
\]
Indeed the inverse tilt contributes
$e^{-s\sum_iX_i(t)}$, and the factor
$e^{-sU_\zeta}$ in $Z(t)$ supplies the missing noise sum.
On $\mathcal A$, $0\le\widetilde W_m-y\le1$, so the
bracketed test lies in $[0,1]$. Replacing its smoothed tilted
product law by the reference law costs at most $mL^{-J}$.
The reference expectation is the same $p_L$ for every $t$.
The remaining prefactor is
$e^{-sy}A(s)^Q(1+o(1))=\mu_L(1+o(1))$.

For a far pair use the joint tilt
$e^{s\sum_i(X_i(t)+X_i(u))}$. In the product $Z(t)Z(u)$ the
two noise integrals are independent copies. After inverse
tilting the test is the product of two tests in $[0,1]$.
The reference one-point paths are independent under the
reference two-point law, so its expectation is exactly
$p_L^2$. It is this reference independence, supplied by
two-point Fourier separation, that is used in
\eqref{eq:far-pair}. Since
\[
 \frac{mL^{-J}}{p_L}\le L^{1-J+C_*},
 \qquad
 \frac{mL^{-J}}{p_L^2}\le L^{1-J+2C_*},
\]
both errors tend to zero under the one chosen inequality
$J>2C_*+12$. Relative normalizer errors also tend to zero.
No lower bound on the unnormalized rare probability
$\mu_Lp_L$ is used for the smoothing step; its exponential
scale has already been extracted exactly.
% END EXPANSION first_far_moments


For non-far pairs let class $k$ be
\[
 e^{-a_k+\Delta}\le d_R(t,u)<e^{-a_{k-1}+\Delta},\quad1\le k\le m,
\]
including every smaller distance in class $m$. The measure of that
class in $\T^2$ is at most
\begin{equation}\label{eq:near-area}
 CR^2 e^{-a_{k-1}+\Delta}.
\end{equation}
Every nonzero integer linear map from $\T^2$ to $\T$ preserves uniform
measure, so a union bound over the $O(R^2)$ maps proves this estimate.
Blocks after $k$ are two-point separated.

On both noisy path events, the actual terminal sums are at least
$y-m\delta$, whereas the prefix at $u$ through block $k$ is at most
$B_k+m\delta$. Applying exponential Markov to the terminals and using
that prefix bound yields
\begin{align}\label{eq:near-moment}
 \E_H[Z(t)Z(u)]
 &\le e^{-2sy+sB_k+O(sm\delta)}
    \E_H e^{s\sum_{i\le k}X_i(t)}
    \E_H e^{s\sum_{i>k}(X_i(t)+X_i(u))}\\
 &\le \exp\{-2sy+sB_k+\lambda S_k+2\lambda(Q-S_k)+o(1)\}.
 \nonumber
\end{align}
The two factors in the first line use disjoint blocks, which are product
kernels under $\E_H$. The prefix uses only one-point separation of
$t\in\mathcal D_L$; it does not assume independence of the two
unseparated prefixes. The suffix uses two-point separation after $k$.
For $k=m$ it is exactly one. The bounded noise weights cause only the
displayed $O(sm\delta)$ error.
Since $sB_k=a_k+\lambda S_k-sG$, division by $\mu_L^2p_L^2$ gives
\begin{equation}\label{eq:near-normalized}
 \frac{\E_H[Z(t)Z(u)]}{\mu_L^2p_L^2}
 \le C e^{a_k}L^{-s_{\min}g_*+2C_*}.
\end{equation}
Multiply by \eqref{eq:near-area} and sum over $m\le L$ classes.
Because $a_k-a_{k-1}=\omega\le A_0\ell$, the normalized near contribution
is at most
\begin{equation}\label{eq:near-sum}
 C L^{1+2u_2+u_3+A_0-s_{\min}g_*+2C_*}=o(1).
\end{equation}
Thus the exponential correlation bound and the arithmetic area cancel
at the correct scale. The polynomial gap $G$ beats all remaining errors.

The total measure of non-far pairs is $O(R^2e^{-r+\Delta})=o(1)$.
Equations~\eqref{eq:first-moment}, \eqref{eq:far-pair}, and
\eqref{eq:near-sum} show
\[
 \frac{\E_H Z_{\mathcal D}^2}{(\E_H Z_{\mathcal D})^2}\longrightarrow1
\]
uniformly over the stated environments and all choices of $\mathcal D$.
Chebyshev's inequality proves
\begin{equation}\label{eq:middle-lower}
 \Pp_H\left(\sup_{t\in\mathcal D}X_{\rm mid}(t)
             <\frac{a_*-r+\lambda Q}{s}-m\delta\right)=o(1).
\end{equation}
All integrals defining these functionals are justified by nonnegativity;
$Z(t)\le e^{sm\delta}$. Singularities of the logarithm at roots have zero
tilted mass and cannot create a high point.

% BEGIN EXPANSION near_and_variance
\paragraph{The last near class and the uniform second-moment conclusion.}
The inequality for near pairs can be verified before taking
any expectation. For fixed noises write $T_t,T_u$ for the
two actual terminal heights, $P_t,P_u$ for the actual prefixes
through block $k$, and $R_t,R_u$ for the two actual
suffix heights after block $k$. On the two path events,
$T_t,T_u\ge y-m\delta$ and $P_u\le B_k+m\delta$.
After dropping the indicators and bounding the two noise
weights, the integrand is bounded by
\[
 \exp\{-2sy+sB_k+Csm\delta\}
       \exp\{sP_t+s(R_t+R_u)\}.
\]
This follows by multiplying the two terminal exponential
Markov bounds and then using $T_u=P_u+R_u$.
Only the $t$-prefix remains random in the bound; the
$u$-prefix has been bounded by its barrier. Products across
prefix and suffix blocks factor under the actual harmonic
product law. At $k=m$ both suffixes are empty, their
exponential factor and expectation are one, and the same
bound is valid. Thus pairs with arbitrarily small or zero
arithmetic separation are included in the last class.

For the area estimate, a nonzero integer vector $(j,l)$
defines a surjective homomorphism
$(t,u)\mapsto jt+lu$ of tori, so the inverse image of a
distance-$\epsilon$ interval has Haar measure at most
$2\epsilon$. Summing over at most $(2R+1)^2-1$ vectors
proves \eqref{eq:near-area}, also if some vectors have
common divisors or one coordinate is zero. The cancellation
in the normalized moment is exact:
\[
 -2sy+sB_k+\lambda S_k+2\lambda(Q-S_k)
       -\log(\mu_L^2)=a_k-sG.
\]
Hence its only nonpolynomial factor, $e^{a_k}$, is canceled
by $e^{-a_{k-1}}$ from area, leaving
$e^\omega\le L^{A_0}$. The finite last class uses the same
area bound and requires no separated suffix.

For completeness, the passage from these bounds to probability
is uniform in the measurable set $\mathcal D$. Put
$d_{\mathcal D}=\Leb(\mathcal D)\ge1/2$ and
$a_L=d_{\mathcal D}\mu_Lp_L>0$. The first moment is
$a_L(1+o(1))$. On far pairs the second-moment integrand is
at most $(1+o(1))\mu_L^2p_L^2$, whose integral over
$\mathcal D^2$ is at most $(1+o(1))a_L^2$. The entire
near integral is $o(\mu_L^2p_L^2)=o(a_L^2)$, since
$d_{\mathcal D}^{-2}\le4$. Thus
\[
 \Pp_H(Z_{\mathcal D}=0)
 \le\frac{\Var_H(Z_{\mathcal D})}{(\E_H Z_{\mathcal D})^2}
 =\frac{\E_H Z_{\mathcal D}^2}{(\E_H Z_{\mathcal D})^2}-1
 =o(1).
\]
For example, if the first moment is at least
$(1-\epsilon)a_L$ and the second is at most
$(1+\epsilon)a_L^2$, with $0<\epsilon<1$, the last bound
is at most $(1+\epsilon)/(1-\epsilon)^2-1$.
All bounds depend on $\mathcal D$ only through its lower
mass bound. Also $Z_{\mathcal D}>0$ entails a point and
a noise realization with terminal height at least $y$;
the actual height there is at least $y-m\delta$.
This completes the lower estimate with its required
uniformity over every such $\mathcal D$.
% END EXPANSION near_and_variance


\subsection{The upper kernel calculation}
For a nonzero polynomial $P$ of degree $d\ge1$,
\begin{equation}\label{eq:poly-estimates}
 \sup_{|z|\le1}|P'(z)|\le ed\|P\|_\infty,\qquad
 \int_\T |P(e(t))|^s\dd t\ge\frac{c_s}{d}\|P\|_\infty^s.
\end{equation}
For the first estimate, the maximum principle applied to $P(z)/z^d$
outside the disk gives $|P(z)|\le|z|^d\|P\|_\infty$. Cauchy's estimate
on a circle of radius $1/d$ about a point in the unit disk then gives the
derivative bound. An arc of length $c/d$ around a boundary maximizer
has modulus at least half the maximum, proving the second estimate.
The constants may be uniform for $s$ in its compact interval.
The middle polynomial has degree at most $Qb\le CLb$ deterministically.
Hence Lemma~\ref{lem:integrated} and Markov's inequality imply, for a
sufficiently large fixed $C_1$,
\begin{equation}\label{eq:middle-upper}
 \Pp_H\left(\max_t X_{\rm mid}(t)
       >\frac{a_*+\lambda Q}{s}+C_1\ell\right)
 \le C L^{1-s_{\min}C_1}=o(1).
\end{equation}
Indeed this event forces the integral in \eqref{eq:integrated} to be at
least $cA(s)^Q L^{sC_1-1}$. No spatial discretization assumption is
needed.

\subsection{Restoring low cycles before the coefficient is completed}
Let $q_0\le Cr$ and write $X_0$ for the low field. For every deterministic
low configuration,
\[
 \int_\T(-X_0(t))_+\dd t\le c_Vq_0,\qquad
 c_V=\int_\T(-\log|1-e(t)|)_+\dd t<\infty.
\]
Choose fixed $C_2$ so large that $\{X_0<-C_2r\}$ has measure at most
$1/8$. For large $n$ its complement intersected with $\mathcal D_L$
is a measurable set $\mathcal D$ of measure at least $1/2$.
Under the harmonic product measure it depends only on low frequencies,
which are independent of the middle frequencies. The bound
\eqref{eq:middle-lower} is uniform over every such measurable set, so it
applies conditionally to this choice. Together with
$X_0(t)\le q_0\log2$, \eqref{eq:middle-upper}, and
$q_0=O(r)$, it follows that
\begin{equation}\label{eq:short-H}
 \Pp_H\left(\left|\log\left\|\prod_{j\le b}(1-z^j)^{C_j}\right\|_\infty
             -\frac{a_*+\lambda(k-l)}s\right|>C_3\ell\right)=o(1).
\end{equation}
Here the counts are fixed, $q_0+Q=k-l$, $r=r_*\ell$, and the constants
remain uniform on all regular count vectors with the bounds above.

Now apply Theorem~\ref{thm:positive-transfer} to the full short-cycle
configuration, including the low part, and to the indicator of the event
in \eqref{eq:short-H}. Its degree is at most $kb$ and its coefficients
are positive. The actual conditional failure probability, given the same
count vector, is at most $(1+O(\eps_n))$ times its harmonic-product
failure probability. This is a single comparison after all projections;
there is no factor $(1+\eps_n)^m$ and no subtraction of an error bigger
than the high-point probability. Finally Lemma~\ref{lem:regular} shows
that these count vectors have actual conditional probability tending to
one. Therefore \eqref{eq:short-H} holds under $\Pp_{n,k}$ as well,
uniformly in $k$, with its random $l=L_*\le C\ell$.

% BEGIN EXPANSION low_and_transfer
\paragraph{Why the random low-frequency set is admissible.}
The deterministic estimate for $X_0$ uses
$(-\sum_vx_v)_+\le\sum_v(-x_v)_+$ and invariance of Haar
measure under $t\mapsto jt$, for each positive integer
frequency $j$. Consequently
\[
 \Leb\{t:X_0(t)<-C_2r\}\le\frac{c_Vq_0}{C_2r}\le\frac18
\]
after one fixed choice of $C_2$. Roots form a finite set
and do not affect this integral or the measure bound.
Together with $\Leb(\mathcal D_L^c)=o(1)$ this gives the
claimed measurable $\mathcal D$ with mass at least one half.
For every realized low configuration apply the just-proved
uniform lower estimate to that particular set. The middle
harmonic sample law is unchanged by this conditioning.
Integrating over low configurations therefore preserves
the same upper bound on its failure probability.

The centering adjustment at this step is explicit:
\[
 \frac{a_*+\lambda(q_0+Q)}s
       -\left(\frac{a_*-r+\lambda Q}s-C_2r-m\delta\right)
 =\frac{r+\lambda q_0}s+C_2r+m\delta=O(\ell).
\]
For the upper estimate one adds at most $q_0\log2$
to the middle maximum. The compact bounds on $s$ and
$\lambda(s)$, $q_0\le Cr$, and the fixed value of $r_*$
make these constants uniform.

Finally, let $\mathcal R_{n,k}$ denote the set of count
vectors obeying all regularity and reservoir-count bounds.
Its actual conditional probability tends uniformly to one.
On this set the normalized order estimate bounds the
conditional short-field failure by a common quantity
$\frac{1+\eps_n}{1-\eps_n}o(1)$, independent of the
individual count vector. Summing these conditional
probabilities bounds the unconditional failure by
\[
 \Pp_{n,k}(\mathcal R_{n,k}^{\,c})
       +\frac{1+\eps_n}{1-\eps_n}o(1)=o(1).
\]
Only count vectors of positive actual mass enter this
sum. The positive transfer also supplies positivity
of the denominator whenever the reference mass is
positive in the reservoir range.
% END EXPANSION low_and_transfer


\subsection{Long-cycle insertion with an exact conditional moment}
\begin{lemma}[Insertion stability]\label{lem:insertion}
Let $P$ be nonzero of degree $d\ge1$ and
$Q(z)=\prod_v(1-z^{j_v})$, where all $j_v$ are positive integers. Then
\begin{equation}\label{eq:insertion}
 -\log2-2ed\sum_vj_v^{-1}
 \le\log\|PQ\|_\infty-\log\|P\|_\infty
 \le(\#v)\log2.
\end{equation}
For $d=0$ the lower bound can be replaced by zero.
\end{lemma}
\begin{proof}
The upper bound is pointwise. At a maximizer $z_*$ of $P$ on the unit
circle take $\rho=1-(2ed)^{-1}$. Equation~\eqref{eq:poly-estimates}
gives $|P(\rho z_*)|\ge\|P\|_\infty/2$. Also
\[
 |Q(\rho z_*)|\ge\prod_v(1-\rho^{j_v}),\qquad
 -\log(1-e^{-x})\le(e^x-1)^{-1}\le x^{-1}\quad(x>0).
\]
Since $\rho^j\le e^{-j/(2ed)}$, the claimed lower bound follows from
the maximum modulus principle for $PQ$. When $P$ is constant use
$Q(0)=1$.
\end{proof}

\begin{lemma}[Conditional cross-mass bound]\label{lem:crossmass}
For $k/L$ in the fixed compact interval and $b\le n/4$,
\begin{equation}\label{eq:crossmass}
 \E_{n,k}\left[
 \left(\sum_{j\le b}jC_j\right)
 \left(\sum_{l>b}\frac{C_l}{l}\right)\right]\le C
\end{equation}
for all sufficiently large $n$, uniformly in $b,k$.
\end{lemma}
\begin{proof}
Deleting one $j$-cycle and one $l$-cycle in \eqref{eq:conditional-law},
where $j\ne l$, gives the exact identity
\[
 \E_{n,k}(C_jC_l)=\frac1{jl}\frac{a_{n-j-l,k-2}}{a_{n,k}}.
\]
Thus the expectation in \eqref{eq:crossmass} equals
\begin{equation}\label{eq:cross-exact}
 \sum_{\substack{j\le b<l\\j+l\le n}}
 \frac1{l^2}\frac{a_{n-j-l,k-2}}{a_{n,k}}.
\end{equation}
If $l\le n/2$, then $n-j-l\ge n/4$. Formula~\eqref{eq:akn}, applied
at sizes in $[n/4,n]$, bounds the coefficient ratio by a constant
uniformly: $\log(n-j-l)=L+O(1)$ and
$k(k-1)/L^2$ is bounded. This part of the sum is at most
$Cb\sum_{l>b}l^{-2}\le C$.

For $l>n/2$, use $l^{-2}\le4/n^2$ and the exact cumulative identity
\begin{equation}\label{eq:sum-a}
 \sum_{r=0}^n a_{r,k-2}
 =[u^{k-2}]h_n(u+1)
 =n a_{n,k-1}+a_{n,k-2}.
\end{equation}
It follows from $h_n(u+1)=(n+u)h_n(u)/u$.
The ratios of the two coefficients on the right to $a_{n,k}$ are
bounded by \eqref{eq:akn}. Hence this part of \eqref{eq:cross-exact}
is at most $Cb/n$. This proves the assertion. The relevant $k$ exceed
two for large $n$, so all displayed coefficients are well defined with
the conventions in \eqref{eq:basic-coeff}.
\end{proof}

\begin{proof}[Completion of Theorem~\ref{thm:main}]
Under $\Pp_{n,k}$, \eqref{eq:short-H} holds with probability tending to
one, uniformly, and $L_*\le C\ell$. For the upper bound multiply back
the $L_*$ long factors, each of supremum at most two. Since
$L-a_*=4\ell+o(1)$, this changes the comparison with
$(L+\lambda k)/s$ by at most a fixed multiple of $\ell$.

For the lower bound apply Lemma~\ref{lem:insertion} with
$d=\sum_{j\le b}jC_j$. By Lemma~\ref{lem:crossmass} and Markov's
inequality, $2ed\sum_{l>b}C_l/l\le\ell$ with probability tending to
one uniformly. Thus reinsertion decreases the maximum by at most
$\ell+\log2$ on that event. The difference between the two centers is
\[
 \frac{L+\lambda k}{s}-\frac{a_*+\lambda(k-L_*)}{s}
   =\frac{L-a_*+\lambda L_*}{s}=O(\ell).
\]
Combining with \eqref{eq:short-H} proves the lower half of
\eqref{eq:main}, and hence the theorem. No lower-bound estimate for a
Poisson spectral maximum was used in this completion.

% BEGIN EXPANSION fixed_constant_completion
\paragraph{A fixed localization constant before the error tolerance.}
Here is an explicit final choice to keep track of the
quantifiers. Let $C_{\rm sh}$ be the fixed constant in
the short-field estimate, and let $C_{\rm long}$ be
fixed so that $L_*\le C_{\rm long}\ell$ with uniformly
high probability. Write
$\lambda_{\max}=\max_{\kappa}\lambda(s_\kappa)$.
The identity $\lambda(0)=\lambda'(0)=0$ and
$\lambda''>0$ gives $\lambda_{\max}>0$.
For all sufficiently large $n$,
$0\le L-a_*\le5\ell$, so on this count event the
difference of centers is nonnegative and at most
$C_{\rm ctr}\ell$, where
\[
 C_{\rm ctr}=\frac{5+\lambda_{\max}C_{\rm long}}{s_{\min}},
 \qquad
 C=C_{\rm sh}+C_{\rm ctr}
           +C_{\rm long}\log2+1+\log2.
\]
Increase the threshold for $n$ so that $\ell\ge1$.
On the insertion event
$2ed\sum_{l>b}C_l/l\le\ell$, the lower loss is
at most $(1+\log2)\ell$; the upper loss is at most
$C_{\rm long}\ell\log2$. Both sides of the desired
inequality then follow with the displayed $C$.
If $d=0$, the constant-polynomial version of insertion
gives an even smaller lower loss.

The complement of this intersection has probability
at most the sum of the short-field failure, the
long-count failure, and
\[
 \Pp_{n,k}\!\left(2ed\sum_{l>b}C_l/l>\ell\right)
 \le\frac{2e}{\ell}
       \E_{n,k}\!\left[d\sum_{l>b}C_l/l\right]
 \le C_{\rm ins}/\ell.
\]
Each tends uniformly to zero. In particular, for any
$\epsilon>0$ one subsequently chooses a threshold
making these three terms at most $\epsilon/3$ each.
The constants defining $C$ have already been fixed.
The integer window in \eqref{eq:main} is nonempty
for large $n$ because its length tends to infinity,
and all its integers lie between $1$ and $n$.
For each such integer $k$, a permutation with one
cycle of length $n-k+1$ and $k-1$ fixed points shows
that $a_{n,k}>0$. Thus the conditional laws and the
supremum in the theorem have their literal meaning.
% END EXPANSION fixed_constant_completion

\end{proof}

\section{Consequences and the limits of the conclusion}\label{sec:consequences}
For fixed $\theta>0$ let $K_n$ have its ordinary Ewens law. Its exact
probability generating function is $h_n(\theta u)/h_n(\theta)$.
The gamma-ratio estimate gives
\[
 \frac{K_n-\theta\log n}{\sqrt{\theta\log n}}\Longrightarrow G,
 \qquad G\sim\mathcal N(0,1),
\]
and $K_n/\log n\to\theta$ in probability. For example the characteristic
function follows by putting $u=e^{it/\sqrt{\theta\log n}}$ in that
ratio; the exponent tends to $-t^2/2$, and the gamma prefactor tends to
one. Differentiation of \eqref{eq:critical} gives
\begin{equation}\label{eq:v-derivatives}
 v'(\kappa)=\frac{\lambda(s_\kappa)}{s_\kappa},\qquad
 v''(\kappa)=-\frac1{\kappa^3s_\kappa^3\lambda''(s_\kappa)}.
\end{equation}
Uniform exact-cycle localization and a Taylor expansion therefore imply
\begin{corollary}[Random centering in the ordinary Ewens law]
\label{cor:ewens}
For every fixed $\theta>0$, with $a_\theta=\lambda(s_\theta)/s_\theta$,
\[
 M_n=v(\theta)\log n+a_\theta(K_n-\theta\log n)+\Op(\log\log n).
\]
In fact a fixed-constant $\log\log n$ bound holds with probability tending
to one. Jointly,
\[
 \left(\frac{K_n-\theta\log n}{\sqrt{\theta\log n}},
       \frac{M_n-v(\theta)\log n}{a_\theta\sqrt{\theta\log n}}\right)
 \Longrightarrow(G,G).
\]
\end{corollary}
\begin{proof}
Use Theorem~\ref{thm:main} on a compact interval containing $\theta$ in
its interior. The quadratic Taylor remainder in
$L v(K_n/L)$ is $O_{\Pp}(1)$, since
$K_n-\theta L=O_{\Pp}(\sqrt L)$. This is smaller than $\ell\to\infty$;
the stated fixed-constant bound and the joint limit follow.

% BEGIN EXPANSION ordinary_center
\paragraph{The Taylor remainder and the common Gaussian coordinate.}
To detail the deterministic expansion, implicit differentiation
in \eqref{eq:critical} yields
\[
 \frac{\dd s_\kappa}{\dd\kappa}
       =-\frac1{\kappa^2s_\kappa\lambda''(s_\kappa)}.
\]
Substitution into $v(\kappa)=\kappa\lambda'(s_\kappa)$
gives the two derivatives in \eqref{eq:v-derivatives}.
In particular $v''$ is bounded on a sufficiently small
fixed compact neighborhood of $\theta$. Since
$\lambda(s)>0$ for $s>0$, the denominator
$a_\theta\sqrt{\theta L}$ in the corollary is positive.
For $k/L$ in that neighborhood Taylor's formula gives
the deterministic bound
\[
 \left|L v(k/L)-L v(\theta)
              -a_\theta(k-\theta L)\right|
       \le C_\theta\frac{(k-\theta L)^2}{L}.
\]
The standardized cycle counts are tight by their CLT,
so the right side, at $k=K_n$, is $O_{\Pp}(1)$.
The probability of leaving the neighborhood tends
to zero by the count law of large numbers.

Theorem~\ref{thm:main}, averaged over $K_n$ inside
this neighborhood, gives a fixed $C_0$ for which
$\Pp(|M_n-Lv(K_n/L)|>C_0\ell)\to0$.
For any family $R_n=O_{\Pp}(1)$ and any deterministic
$\ell\to\infty$, $\Pp(|R_n|>\ell)\to0$: given an
error tolerance, first choose a uniform tightness
bound and then let $\ell$ exceed it. Applying this
observation to the Taylor remainder proves the
corollary's fixed-constant assertion with $C_0+1$.

Set $Y_n=(K_n-\theta L)/\sqrt{\theta L}$ and let
$V_n$ denote the second coordinate in the corollary.
The bound just proved implies
$V_n-Y_n\to0$ in probability, since
$\ell/\sqrt L\to0$. The continuous map
$x\mapsto(x,x)$ sends the count CLT to
$(Y_n,Y_n)\Rightarrow(G,G)$; the vanishing difference
$(0,V_n-Y_n)$ then gives the displayed joint weak limit.
This is a statement about joint distributions.
It does not require, or assert, convergence of second
moments or a limit for the Pearson correlation.
% END EXPANSION ordinary_center

\end{proof}

\paragraph{The original two-matching encoding.}
If $a$ and $b$ are fixed-point-free involutions on $2n$ labels, with $a$
fixed and $b=gag^{-1}$ uniformly conjugated, the alternating component
counts $C_j$ have the Ewens$(1/2)$ law in the encoding of \cite{Lu}.
On an alternating component of size $2j$, the product $ab$ consists of
two $j$-cycles. Hence
\[
 \det(I-zP_{ab})=\prod_j(1-z^j)^{2C_j}.
\]
If the number of alternating components is fixed to $k$,
Theorem~\ref{thm:main} applies directly: the logarithmic maximum is within
$C\log\log n$ of $2m_{n,k}$, uniformly for $k/\log n$ in a positive
compact interval. This is an exact application of the encoding, not an
additional extremal theorem proved by a different approximation.

% BEGIN EXPANSION matching_encoding
\paragraph{Counting and identifying the actual matching model.}
For completeness, the encoding can be described directly.
Relabel the fixed matching $a$ so that its edges are
the pairs $(i,0),(i,1)$, $1\le i\le n$. For another
matching $b$, the graph with edges of both matchings
has even alternating components. A common edge
$a=b$ is a component on two vertices; it is included.
There are exactly two binary colorings of each
component in which both kinds of edges join opposite
colors. Consequently a matching with $k$ components
admits $2^k$ such colorings.

Given a coloring, denote by $\operatorname{out}(i)$
and $\operatorname{in}(i)$ the two vertices of the
$i$th fixed edge, according to their colors.
The matching $b$ sends $\operatorname{out}(i)$ to
$\operatorname{in}(\pi(i))$ for a unique permutation
$\pi$ of the $n$ fixed edges. Conversely, a choice
of the $n$ orientations and a permutation $\pi$
specifies all these $b$-edges and the coloring.
The alternating components correspond exactly to
the cycles of $\pi$, with twice their vertex counts.
The number of permutations with cycle profile $c$
is $n!\prod_j(j^{c_j}c_j!)^{-1}$. Counting the
colored matchings in these two ways therefore gives
\[
 \#\{b:\text{component profile }c\}
 =2^{\,n-k}n!\prod_j\frac{j^{-c_j}}{c_j!},
 \qquad \sum_jjc_j=n,\quad \sum_jc_j=k.
\]
Dividing these weights by their sum identifies the
component profile as Ewens$(1/2)$, and conditioning
on its component count cancels the factor $(1/2)^k$.
The normalization agrees with
$\#\{b\}=(2n-1)!!=2^n n!h_n(1/2)$.
Uniform conjugation indeed produces the uniform
matching: all fixed-point-free involutions are
conjugate, and every fiber of $g\mapsto gag^{-1}$
is a coset of the same centralizer. Relabeling back
to any fixed $a$ preserves the component profile.

The same description determines the permutation
$ab$. On the out-vertices it acts as $\pi$ and on
the in-vertices as $\pi^{-1}$. A component of size
$2j$ therefore contributes two $j$-cycles, including
two fixed points when $j=1$. A $j$-cycle contributes
$1-z^j$ to $\det(I-zP)$, so the determinant formula
above holds for the actual product permutation.
Writing $P_c(z)=\prod_j(1-z^j)^{c_j}$, one has
$\|P_c^2\|_\infty=\|P_c\|_\infty^2$ and therefore
the actual logarithmic maximum equals
$2\log\|P_c\|_\infty$. Its conditioning variable
is the number $k$ of graph components, whereas
the permutation $ab$ itself has $2k$ cycles.
Applying Theorem~\ref{thm:main} to the size-$n$
profile and multiplying its error bound by two
proves exactly the asserted center $2m_{n,k}$.
% END EXPANSION matching_encoding


\paragraph{What is and is not being asserted.}
The proof above closes the conditional high-probability lower bound
through a coefficient comparison valid for the entire path restriction.
It also supplies the matching upper localization. It does not identify
$M_n-m_{n,K_n}$ after subtracting a particular multiple of $\log\log n$,
nor its distribution on a smaller scale, nor an extremal point process.
A previous fixed-Diophantine moment theorem and a microscopic collision
kernel are not inputs here: reorganizing the generating function into a
reservoir and a positive constrained polynomial avoids needing a uniform
extension of those theorems over all moving singularity clusters.

The proof still uses a final second-moment inequality and elementary
walk estimates. Its coefficient content is not a claim that probability
can be eliminated from a probabilistic problem. The methodological point
is that exact conditioning and all growing-dimensional path restrictions
are handled by proved coefficient and convolution estimates, rather than
being removed by an absolute-error approximation or inserted as a
hypothesis.

\section*{Tools and formal verification}
The mathematical development of this paper was carried out through
iterative work with AI assistants, which also translated the argument
into Lean~4 and ran the verification. The formalization is over the
literal finite permutation model: it covers the conditional law
\eqref{eq:conditional-law}, the characteristic polynomial as the
determinant of the actual permutation matrix, Theorem~\ref{thm:main},
Corollary~\ref{cor:ewens} and the two-matching application of
Section~\ref{sec:consequences}. It was checked with \texttt{trust=0} and
\texttt{debug.skipKernelTC=false}, contains no unproved placeholders,
and uses only the foundational axioms \texttt{propext},
\texttt{Classical.choice} and \texttt{Quot.sound}. The Lean sources, the
locked dependencies and the verification script are available at
\begin{center}
\url{https://github.com/Lzp88/conditional-spectral-extremes}.
\end{center}
Kernel checking establishes the formal statements, not their equivalence
with the prose above, and bears on neither novelty nor priority. The
author is responsible for the mathematical claims.

\begin{thebibliography}{9}
\bibitem{Lu}
Z. Lu, \emph{On counting double centralizers of symmetric groups},
Electronic Journal of Combinatorics \textbf{30}(2) (2023), P2.47.
\bibitem{FFGPP}
P. Flajolet, E. Fusy, X. Gourdon, D. Panario, and N. Pouyanne,
\emph{A hybrid of Darboux's method and singularity analysis in
combinatorial asymptotics}, Electronic Journal of Combinatorics
\textbf{13} (2006), R103. doi:10.37236/1129.
\bibitem{CZ}
N. Cook and O. Zeitouni, \emph{Maximum of the characteristic polynomial
for a random permutation matrix}, Communications on Pure and Applied
Mathematics \textbf{73} (2020), 1660--1731. doi:10.1002/cpa.21899.
\bibitem{HNNZ}
C. Hughes, J. Najnudel, A. Nikeghbali, and D. Zeindler,
\emph{Random permutation matrices under the generalized Ewens measure},
Annals of Applied Probability \textbf{23}(3) (2013), 987--1024.
doi:10.1214/12-AAP862.
\bibitem{DZ}
K. Dang and D. Zeindler, \emph{The characteristic polynomial of a random
permutation matrix at different points}, Stochastic Processes and their
Applications \textbf{124}(1) (2014), 411--439.
doi:10.1016/j.spa.2013.08.003.
\bibitem{Francois}
Q. Fran\c{c}ois, \emph{Characteristic polynomial of generalized Ewens
random permutations}, arXiv:2504.01484v2 (2025).
\bibitem{KMS}
Z. Kabluchko, A. Marynych, and H. Sulzbach,
\emph{Mode and Edgeworth expansion for the Ewens distribution and the
Stirling numbers}, Journal of Integer Sequences \textbf{19} (2016),
Article 16.8.8.
\bibitem{MM}
B. Mallein and P. Mi\l o\'s, \emph{Maximal displacement of a supercritical
branching random walk in a time-inhomogeneous random environment},
Stochastic Processes and their Applications \textbf{129}(9) (2019),
3239--3260. doi:10.1016/j.spa.2018.09.008.
\bibitem{CG}
N. A. Cook and H. Gu, \emph{On the maximum of Poissonian log-correlated
fields}, announced work, listed as in preparation on H. Gu's
\href{https://sites.google.com/view/haotian-gu/research}{research page}
(accessed September 10, 2026). See also H. Gu's
\href{https://calendar.usc.edu/event/probability-and-statistics-seminar-maximum-of-poissonian-log-correlated-fields}{seminar announcement},
University of Southern California, January 14, 2026; the
\href{https://erdoscenter.renyi.hu/sites/default/files/events/documents/short_talk_book_of_abstracts.pdf}{School on Disordered Media abstracts},
Erd\H{o}s Center, January 20--24, 2025, p.~6; and N.~Cook's
\href{https://www-math.umd.edu/gcal_rss.php?html=\&seminar_key=PROB\&year=2026}{probability seminar},
University of Maryland, October 22, 2025.
\end{thebibliography}
\end{document}
